\documentclass[11pt,a4paper]{article}

\usepackage[T1]{fontenc}
\usepackage{mathpazo}

\usepackage{amsmath,amssymb,amsthm,mathtools}
\usepackage{mathrsfs}
\usepackage{geometry}
\usepackage{microtype}
\usepackage{setspace}
\usepackage{titlesec}
\usepackage{titletoc}
\usepackage{fancyhdr}
\usepackage{xcolor}
\usepackage{booktabs}
\usepackage{array}
\usepackage{tikz}
\usetikzlibrary{decorations.pathmorphing,arrows.meta,positioning}
\usepackage{tikz-cd}
\usepackage[hidelinks]{hyperref}

\geometry{
  a4paper,
  top=1.1in,
  bottom=1.1in,
  left=1.25in,
  right=1.25in,
  headheight=15pt
}

\onehalfspacing

\definecolor{rulecolor}{gray}{0.4}

\titleformat{\section}
  {\large\bfseries}
  {\thesection.}
  {0.6em}
  {}
  [\vspace{0.1em}\textcolor{rulecolor}{\titlerule[0.4pt]}]

\titleformat{\subsection}
  {\normalsize\bfseries\itshape}
  {\normalfont\normalsize\bfseries\thesubsection.}
  {0.5em}
  {}

\titlespacing*{\section}{0pt}{1.6em}{0.7em}
\titlespacing*{\subsection}{0pt}{1.1em}{0.4em}

\pagestyle{fancy}
\fancyhf{}
\fancyhead[L]{\small\itshape Constraint Geometry Across Scales}
\fancyhead[R]{\small\itshape Flyxion}
\fancyfoot[C]{\small\thepage}
\renewcommand{\headrulewidth}{0.3pt}

\theoremstyle{plain}
\newtheorem{theorem}{Theorem}[section]
\newtheorem{proposition}[theorem]{Proposition}
\newtheorem{corollary}[theorem]{Corollary}
\theoremstyle{definition}
\newtheorem{definition}[theorem]{Definition}
\theoremstyle{remark}
\newtheorem{remark}[theorem]{Remark}

% ---- Macros ----
\newcommand{\A}{\mathcal{A}}
\newcommand{\C}{\mathcal{C}}
\newcommand{\fM}{\mathfrak{M}}
\newcommand{\Pp}{\mathcal{P}}
\newcommand{\Rr}{\mathcal{R}}
\newcommand{\Svalid}{\mathcal{S}_{\mathrm{valid}}}
\newcommand{\Adm}{\mathcal{A}_{\mathrm{dm}}}
\newcommand{\KES}{\mathfrak{S}}
\newcommand{\RSVP}{\textup{\textsc{rsvp}}}
\newcommand{\CLIO}{\textup{\textsc{clio}}}
\newcommand{\TARTAN}{\textup{\textsc{tartan}}}
\newcommand{\FSA}{\textup{\textsc{fsa}}}
\newcommand{\Omg}{\Omega}
\newcommand{\scrA}{\mathscr{A}}

\setlength{\parindent}{1.5em}
\setlength{\parskip}{0.2em}

\hyphenation{ad-mis-si-bil-i-ty-pre-serv-ing ad-mis-si-bil-i-ty}

% -------------------------------------------------------
\begin{document}

\begin{titlepage}
\vspace*{2cm}
\begin{center}

{\LARGE\bfseries Constraint Geometry Across Scales:\\[0.3em]
Field-Space Autoencoders, Spin--Valley Modes,\\[0.3em]
and the RSVP\,/\,CLIO\,/\,TARTAN Research Program}

\vspace{1.2em}

{\large\itshape Convergent Structure in Climate Compression,\\
Condensed Matter Physics, and Constraint-First Field Theory}

\vspace{2em}

{\large Flyxion}\\[0.3em]
{\small Independent Researcher}

\vspace{1.5em}

{\small May 2026}

\vspace{3em}

\begin{minipage}{0.82\textwidth}
\small\itshape
\noindent\textbf{Abstract.}
Two recent papers, arriving from entirely different empirical domains,
independently converge on a shared structural principle: stable physical
organization is carried by hidden dynamical manifolds, and observable
outputs are secondary projections of that deeper constraint geometry.
The Field-Space Autoencoder demonstrates this computationally, in climate
emulation: geometry-aware, hierarchically explicit, constraint-preserving
compression dramatically outperforms flat representational approaches.
A contemporaneous study of spin--valley collective modes in twisted
WSe$_2$ moiré superlattices demonstrates it physically: the physically
meaningful excitations of the system propagate through charge-decoupled
hidden modes that are invisible to standard local observables and require
direct space-time imaging to detect. This essay examines both convergences
in relation to the \RSVP\ (Relativistic Scalar-Vector Plenum), \CLIO\
(Constraint-Leveraged Inference and Optimization), and \TARTAN\
(Trajectory-Aware Recursive Tiling with Annotated Noise) frameworks ---
not as direct empirical confirmation, but as evidence of a broader
structural pressure forcing multiple independent research lines toward a
geometry-first, admissibility-preserving ontology. We develop the formal
correspondences, generalize their consequences, and identify the conceptual
gap that separates operationally sophisticated systems from a fully
field-theoretic account of constraint-geometry propagation.
\end{minipage}

\end{center}
\end{titlepage}

\tableofcontents
\newpage

% =============================================================
\section{Introduction: The Structural Pressure}
\label{sec:intro}

There is a recurring pattern in the history of scientific computation. A
field develops a representational convention that works well within a
particular parameter regime, and then encounters systematic failure as the
physical systems it models become increasingly demanding of geometric
fidelity. The failure is not usually attributable to insufficient data or
insufficient parameter count. It is attributable to a mismatch between the
topology of the representation and the topology of the underlying process.
When that mismatch becomes severe enough, the field eventually reorganizes
around a new representational ontology --- one that takes the geometry of
the process as primary rather than emergent.

\subsection{Representation Failure as Topological Mismatch}

The distinction between numerical error and structural failure is critical.
When a climate model accumulates bias near the poles, or when a quantum
simulation loses long-range coherence, the conventional response is to
increase resolution, add parameters, or improve the loss function. These
interventions address the symptom. The underlying cause is that the
representational substrate --- the grid, the token sequence, the flat
latent space --- is topologically incompatible with the process it is
meant to model. A representational container whose geometry differs from the
physical geometry does not merely introduce approximation error; it
introduces a systematic bias toward trajectories that are admissible in
the representation but inadmissible in the physics. No increase in
parameter count resolves a topological mismatch. Only a change of
representational ontology can.

This distinction --- between insufficient scale and structural incompatibility
--- is the conceptual pivot of the entire essay. The two empirical papers
examined here matter not because they involve large models or impressive
benchmarks, but because they both, independently and from very different
directions, arrive at architectures and measurement strategies whose
design logic is driven by the recognition that the representational
topology must be compatible with the physical topology. That recognition,
when forced by engineering or experimental necessity, is what we call a
structural pressure.

A further structural feature unifies both empirical systems: both involve
long-horizon dynamics where local correctness is insufficient for global
admissibility. A climate model that is accurate at each time step can still
produce a trajectory that is globally inadmissible --- that drifts to
climatologically impossible states --- because local step-by-step
reconstruction does not enforce global constraint coherence. A measurement
strategy that accurately records local charge density at each point in space
can still be entirely blind to the collective mode structure, because the
relevant physics is carried by globally organized hidden fields rather than
local observables. Whenever admissibility is globally path-dependent, short-horizon
optimization fails. This is the deep unifying constraint across climate
systems, quantum collective modes, cognition, and semantic memory.

Climate modelling is currently undergoing exactly this transition. The
dominant approach for decades has been to discretize atmospheric and oceanic
dynamics onto regular grids, apply convolutional or spectral operators to
those grids, and treat the resulting array of values as a compressed
description of the physical state. This approach has been productive, but
it has a structural defect: the sphere is not the plane, and forcing
spherical dynamics into Euclidean representational containers introduces
systematic distortions that accumulate across temporal integration and
spatial downscaling. The distortions are not merely numerical. They reflect
a genuine topological incompatibility between the representational substrate
and the physical substrate.

The Field-Space Autoencoder (\FSA) paper, which we treat throughout as our
primary empirical reference, is a response to exactly this failure mode.
Its core architectural innovation is to process climate fields natively on
their natural geometric domains rather than on flattened approximations, to
decompose dynamics into an explicit multi-scale hierarchy rather than
leaving scale structure to be inferred by the network, and to treat
cross-scale operators as structure-preserving morphisms rather than as
unconstrained learned transforms. These are not merely engineering
refinements. They reflect a recognizable shift in representational
ontology: from storage-first to geometry-first, from flat coding to
constraint-preserving projection.

The \RSVP\ (Relativistic Scalar-Vector Plenum), \CLIO\ (Constraint-Leveraged
Inference and Optimization), and \TARTAN\ (Trajectory-Aware Recursive
Tiling with Annotated Noise) frameworks have been developed from a
different direction --- theoretical rather than empirical, field-theoretic
rather than engineering-motivated --- but they have been organized around
precisely the same structural commitments. The convergence is therefore
worth examining carefully. It suggests that these principles are not
artifacts of any particular theoretical programme but rather structural
necessities that emerge independently wherever the representational demand
for geometric fidelity becomes severe enough.

\subsection{Convergent Rediscovery Rather Than Derivation}

The epistemic framing of this essay requires explicit statement, because
the risk of misreading is real. The correspondence between the empirical
architectures examined here and the \RSVP/\CLIO/\TARTAN\ framework is not
a claim of priority, derivation, or retrospective validation. Neither the
\FSA\ paper nor the moiré spectroscopy group arrived at their results by
consulting these frameworks. They arrived there under independent empirical
pressure from the physical systems they were studying. The theoretical
frameworks were developed independently, under different pressures, from a
different direction.

What the convergence demonstrates is that the structural commitments in
question are being driven by something real in the domain rather than by
the preferences of any particular research group. When engineers building
climate emulators and physicists imaging quantum collective modes arrive at
overlapping structural principles without consulting each other or any
shared theoretical source, that overlap is evidence of a structural
necessity. The theoretical frameworks provide a vocabulary for stating that
necessity precisely; they do not constitute its source.

This essay proceeds through five stages. In the first, we examine the
\FSA\ architecture on its own terms and identify the specific ways in which
it departs from flat representational conventions. In the second, we
develop the formal correspondence between the \FSA\ and the \RSVP/\CLIO\
frameworks, arguing that the architecture is an engineering instance of
admissibility-preserving projection. In the third, we generalize the
consequences through \TARTAN, showing that the \FSA\ approach implements
several \TARTAN\ principles in a concrete domain. In the fourth, we
examine the spin--valley moiré paper and argue that it provides a physical
analogue of the same structural principles: hidden collective modes
propagating through constraint manifolds that are invisible to standard
local observables. In the fifth, we draw the unified thesis from both
papers, identify the conceptual gap that separates operationally powerful
systems from a fully field-theoretic account, and develop the sheaf and
cohomological formalization of that gap.

% =============================================================
\section{The Field-Space Autoencoder: Architecture and Failure Modes Addressed}
\label{sec:fsa}

\subsection{The Standard Architecture and Its Defects}

The dominant representational pipeline in data-driven climate modelling
follows a by-now familiar pattern: map the physical field to a fixed
discretization, apply learned operators to that discretization, and decode
the result back to a physical prediction. The appeal of this pipeline is
its generality. Convolutional networks, transformers, and their hybrids can
all be applied within it without modification, and the engineering
infrastructure that has been developed for flat image processing transfers
without fundamental change.

The defects are structural. A climate field is defined on the sphere
$S^2$, and the natural symmetries of that domain --- rotational invariance,
the absence of a canonical pole, the curvature of geodesics --- do not
survive the transition to a regular latitude-longitude grid. Near the poles,
grid cells shrink in area while remaining equal in index space, distorting
both the statistics of the representation and the effective receptive fields
of convolutional operators. Interpolation to a flat grid introduces
artifacts that compound across temporal integration steps. Most
significantly, the representational geometry diverges from the physical
geometry: the allowed transformations in index space are not the allowed
transformations in field space.

The \FSA\ paper identifies this as the central failure and organizes its
architecture around avoiding it. Rather than projecting the physical field
onto a flat representational substrate, it operates natively on the HEALPix
grid~\cite{Gorski2005} --- an equal-area spherical discretization that avoids
polar distortion by treating all spatial locations uniformly. The
architectural consequence is that the network never operates on a
representational object whose topology contradicts the topology of the
process it is trying to model. At 64$\times$ compression, the Field-Space
models achieve reconstruction errors (RMSE $\approx 0.28\,^\circ$C) lower
than the HEALPix convolution baseline achieves at only 16$\times$
compression (RMSE $\approx 0.31\,^\circ$C), representing a fourfold
increase in compression efficiency at equivalent accuracy. The baseline
degrades sharply at extreme compression ratios (256$\times$, 1024$\times$)
where the Field-Space models remain viable.

\subsection{Why Flat Latent Spaces Fail}

It is worth distinguishing three distinct fidelity criteria that are
frequently conflated in representation learning: \textit{reconstruction
fidelity}, \textit{dynamical fidelity}, and \textit{admissibility
fidelity}. A model with high reconstruction fidelity produces outputs that
closely resemble the training data in pixel-level or token-level statistics.
A model with high dynamical fidelity produces outputs whose temporal
evolution is statistically consistent with the underlying physical dynamics.
A model with high admissibility fidelity produces only outputs that are
structurally consistent with the constraint geometry of the domain --- outputs
that lie on or near the manifold of physically realizable states, not merely
near the manifold of statistically frequent ones.

Standard flat latent space architectures optimize for reconstruction fidelity
and, to a lesser extent, dynamical fidelity. They do not optimize for
admissibility fidelity, and they provide no mechanism for enforcing it.
The result is that their outputs can be locally accurate --- close to the
training distribution in the sense of standard loss metrics --- while being
globally inadmissible: on trajectories that no physical system could actually
follow. The failure mode is not large error but wrong error: the model
produces states that look plausible in isolation but are collectively
inconsistent with the admissibility structure of the domain.

The \FSA\ architecture's geometric commitments move it toward admissibility
fidelity without explicitly naming it as such. By operating on the native
geometric domain and enforcing explicit cross-scale compatibility, the
architecture makes it structurally harder to produce globally inadmissible
outputs. The latent space organization that results is evidence that this
structural hardening is working: the compressed representations organize
themselves according to the dynamical geometry of the physical system, not
merely according to reconstruction efficiency.

\subsection{Multi-Scale Residual Decomposition}

The second major architectural innovation is the explicit multi-scale
decomposition strategy. Rather than asking a single network to learn all
spatial scales simultaneously from raw data, the \FSA\ architecture
decomposes the field into a hierarchy of scale-specific representations:
\[
X = X_0 \oplus \bigoplus_{k=1}^{n} R_k,
\]
where $X_0$ is a coarse global structure and each $R_k$ is a residual
correction at progressively finer spatial resolution.

This decomposition is not merely a computational convenience. It encodes a
specific ontological commitment: large-scale structure is primary, and
fine-scale structure exists conditionally upon it. The coarse manifold
$X_0$ acts as an admissibility regulator for the residual layers. A
fine-scale refinement $R_k$ is constrained to be compatible with the
coarser structures at all lower resolutions, not merely with the
immediately preceding layer. This creates a cascade of compatibility
constraints that strongly limits the space of admissible high-resolution
reconstructions.

The architectural specification of the cross-scale operators makes this
explicit. The compression operator takes fields at multiple resolutions as
input and produces representations at an intermediate resolution:
\[
\mathcal{F}_{k \to j}:\{r^{(k)}, r^{(k+1)}, \ldots\} \to \{r^{(j)}\}.
\]
The decompression operator inverts this, reconstructing finer resolutions
from coarser inputs:
\[
\mathcal{F}^{-1}_{j \to k}:\{r^{(j)}, r^{(j-1)}, \ldots\} \to \{r^{(k)}\}.
\]
The explicit specification of these operators as functions of multiple
scale strata, rather than as point-to-point mappings between adjacent
scales, is a strong architectural commitment to the idea that compatibility
is a global rather than local constraint.

\subsection{Latent Space Organization}

The most philosophically significant result in the \FSA\ paper is the
spontaneous organization of the latent space. The authors analyse the
compressed representations of 84 years of daily ERA5 near-surface air
temperature data (1940--2024) using t-SNE dimensionality reduction. The
two Field-Space Autoencoder variants --- convolutional and transformer-based
--- produce latent spaces that are, without any explicit temporal or
climatological supervision, organized into cyclic seasonal trajectories
when coloured by month, and a smooth directional drift from earlier to
more recent states when coloured by year. Both architectures were trained
\textit{exclusively on residual fields}, deliberately excluding the coarse
base component that contains most of the seasonal and warming signal. The
organized latent structure therefore cannot be attributed to exposure to
large-scale climatological information during training; it emerges from the
architecture's representational commitments alone.

In sharp contrast, the HEALPix convolution baseline --- trained on the full
raw fields including the seasonal cycle and warming trend --- produces a
latent space with no comparably coherent structure in the t-SNE projection.
The Field-Space architecture, operating only on fine-scale residuals,
recovers the dynamical geometry of the full climate system. The baseline,
operating on the full field, loses it. This reversal is the empirical result
that most directly confirms the essay's central argument: surface fidelity
and dynamical fidelity are orthogonal properties of a representation, and
optimizing for one can actively impede the other.

\subsection{Compression as Constraint Transport}

The \FSA\ field-space operators do something subtler than dimensionality
reduction. They transport the compatibility structure of the field across
scale boundaries. Each operator $\mathcal{F}_{k \to j}$ takes as input the
field representations at multiple resolution strata and produces an output
that is simultaneously consistent with all of them. This is not compression
in the information-theoretic sense of finding a shorter description of the
same content. It is compression in the constraint-geometric sense of finding
a representation that respects the admissibility conditions imposed by the
coarser scales.

The distinction matters because it determines what generalizes. A system
that compresses by removing information discards the fine-scale structure
and cannot recover it without additional data. A system that compresses by
transporting compatibility structure retains the \textit{conditions} under
which fine-scale structure is admissible, and can therefore reconstruct
compatible fine-scale structure for any input that satisfies those conditions
--- including inputs at resolutions never seen during training. Zero-shot
super-resolution is the empirical signature of constraint transport rather
than information reduction.

\subsection{Zero-Shot Super-Resolution}

The \FSA\ paper demonstrates a capability with direct theoretical
significance. The transformer-based model, trained at a fixed 64$\times$
compression ratio on ERA5 data at HEALPix level~8 (approximately 30~km
resolution), can produce outputs at HEALPix levels 7 and~6 --- corresponding
to 4$\times$ and 16$\times$ super-resolution --- by setting finer residual
levels to zero at inference time. The model has never seen the pairing of
level-6 input with level-8 output during training. The transformer exhibits
only a small decrease in skill at these novel resolution configurations and
no spatial artifacts. The convolutional variant, by contrast, degrades
rapidly and produces prominent HEALPix grid artifacts at level-6 input,
suggesting it overfit to the specific residual resolutions seen during
training rather than learning transferable compatibility conditions.

The paper also applies the pretrained transformer to MPI-ESM1.2-HR
historical simulations natively produced at approximately 100~km
resolution --- a domain the model was never trained on. After encoding at
HEALPix level 6 and decoding to level 8, the model synthesizes
high-frequency spatial structures entirely absent from the source simulations,
as confirmed by spectral analysis showing maintained power at high multipoles
where the climate model falls off. This is sparse reconstruction from
constraint geometry operating across both resolution and model boundaries.
The model has not memorized fine-scale ERA5 patterns; it has learned the
admissibility conditions relating coarse and fine scales, and those conditions
transfer.

\subsection{Feature Learning Versus Constraint Learning}

Most representation-learning systems are implicitly optimised to learn
\textit{features}: statistical regularities present on the observable surface
of the training distribution. A feature is a pattern that recurs across
training examples; feature learning is the process of identifying and
encoding those patterns in a form that supports downstream prediction.
The \FSA\ results suggest a fundamentally different interpretation of what
successful generalisation requires.

The architecture does not succeed because it has memorised high-frequency
climate features. It succeeds because it has learned the \textit{compatibility
relations} that govern which fine-scale structures are admissible given
coarse-scale conditions. These compatibility relations are not surface
patterns; they are structural constraints that hold across domains, resolutions,
and model boundaries because they are properties of the physical system
rather than of the training corpus.

The distinction matters because feature-learning systems and
constraint-learning systems fail differently. A feature-learning system
fails whenever the statistical surface of the input diverges from the training
distribution: when the domain shifts, the features no longer match, and
generalization collapses. A constraint-learning system fails only when the
underlying admissibility structure itself changes --- a much more severe
condition that corresponds to genuine physical change rather than mere
distributional shift. The former is distributionally fragile; the latter is
structurally robust.

The convolutional baseline's behaviour at zero-shot resolution makes this
contrast sharp. It overfits to the specific HEALPix residual resolutions
present during training and produces artifacts when these change, because it
has learned features of the training residual distribution. The transformer
exhibits no such artifacts because it has learned the compatibility conditions
that relate scale strata, and those conditions do not depend on which specific
strata were used during training. This asymmetry is the operational signature
of constraint learning. It also provides a cleaner conceptual target for
the scaling critique developed later: the question is not how many features
a system has memorised but whether it has learned the constraint structure
of the domain.

% =============================================================
\section{The \RSVP\ Framework: Fields, Admissibility, and Trajectory Geometry}
\label{sec:rsvp}

\subsection{Foundational Ontology}

The Relativistic Scalar-Vector Plenum framework takes as its fundamental
ontological commitment the idea that the relevant objects of a physical
theory are not states but \textit{fields over possibility spaces}. A state
is a momentary cross-section of a field; a trajectory is a sequence of
such cross-sections satisfying admissibility constraints; and the
accessible future is the set of trajectories that are reachable given the
current field configuration and accumulated history.

Formally, the \RSVP\ state is a triple
\[
X_t = (\varphi_t, \mathbf{v}_t, S_t),
\]
where $\varphi_t$ is the scalar plenum density governing long-range
coherence, $\mathbf{v}_t$ is the vector flow field governing directional
transport, and $S_t$ is the local entropy density governing the volume of
admissible refinements. The dynamics of each component are coupled:

\begin{align}
\partial_t \varphi &= D_\varphi \nabla^2 \varphi
  - \nabla \cdot (\varphi \mathbf{v})
  + R_\varphi - \lambda_\varphi S, \label{eq:phi}\\
\partial_t \mathbf{v} &= -(\mathbf{v} \cdot \nabla)\mathbf{v}
  + \nu \nabla^2 \mathbf{v}
  - \nabla \varphi
  + \mathcal{T}(\mathbf{v}), \label{eq:v}\\
\partial_t S &= D_S \nabla^2 S
  + \sigma(\varphi, \mathbf{v})
  - \gamma S. \label{eq:S}
\end{align}

The coupling term $\lambda_\varphi S$ in equation~\eqref{eq:phi} is
significant: local entropy density suppresses scalar coherence. Regions
of high entropy are regions where many trajectories remain admissible,
and coherent structure is correspondingly weaker. The scalar field $\varphi$
should therefore be understood not merely as a density field but as a
coherence regulator.

The transport operator $\mathcal{T}(\mathbf{v})$ in equation~\eqref{eq:v}
is equally important. It acts as a geometric connection on the flow field:
it absorbs local topological defects --- vortices, phase slips, singular
points in the vector field --- that would otherwise break the global
coherence of the flow. In the spin--valley context developed later in this
essay, $\mathcal{T}(\mathbf{v})$ plays the role of the gauge-like term
that couples the superfluid velocity to the underlying moiré superlattice
potential, ensuring that the collective flow remains globally admissible
despite local geometric irregularities. This structural symmetry --- between
$\mathcal{T}(\mathbf{v})$ as defect absorber in \RSVP\ and the phase mode
$\mathbf{v}_t \sim \nabla\theta$ as coherence carrier in the spin--valley
superfluid --- is one of the most precise formal correspondences between the
theoretical framework and the condensed matter experiment.

\subsection{Admissibility and the Accessibility Structure}

The central theoretical construct in \RSVP\ is the admissibility function
$\kappa(\omega, H_t)$, which assigns a viability score to each candidate
transition $\omega$ given the accumulated event history $H_t$. The
admissible set at time $t$ is
\[
\Adm(X_t) = \bigl\{\omega \in \Omg_t \mid \kappa(\omega, H_t) \geq \kappa_*\bigr\},
\]
where $\kappa_*$ is the viability threshold below which transitions are
physically or structurally prohibited. The event history $H_t =
\{e_0, e_1, \ldots, e_t\}$ is append-only: transitions from $\Omg_t$
select events into $H_{t+1}$, and this selection is irreversible.

The global model space $\fM = \{s_t \mid t \in \mathbb{R}\}$ supports a
family of history-conditioned accessibility slices
\[
\A_t = \bigl\{s \in \fM \mid s \text{ reachable given } H_t\bigr\}.
\]
Temporal asymmetry is encoded in the strict monotone shrinkage
\[
\A_{t+1}^{\mathrm{consistent}} \subsetneq \A_t^{\mathrm{counterfactual}}:
\]
the consistent slice of states compatible with $H_{t+1}$ is always
strictly smaller than the counterfactual envelope of states that were
accessible before the transition at $t$.

This structure provides a formal account of irreversibility that does not
depend on thermodynamic entropy alone. A transition is irreversible not
merely because it increases entropy but because it eliminates branches from
the accessibility structure that cannot be recovered by any subsequent
sequence of admissible transitions.

\subsection{Projection, Compression, and Admissibility Preservation}

The central epistemological problem in \RSVP\ is what happens when a
field-theoretic system is projected onto a lower-dimensional representational
manifold. Let $\pi: X_t \to M$ denote such a projection, where $X_t$ is
the trajectory space and $M$ is the compressed manifold. The projection is
admissibility-preserving if and only if admissible transitions in $X_t$
map to admissible transitions in $M$:
\[
\omega \in \Adm(X_t) \;\Rightarrow\; \pi(\omega) \in \Adm(M).
\]
When this condition fails, the compressed representation supports
trajectories that would be inadmissible in the underlying physical system.
The model learns to navigate the compressed manifold rather than the
physical one, and generalization breaks down precisely at the boundaries
where the two manifolds diverge.

This failure mode is not merely a theoretical concern. It is the failure
that the \FSA\ paper is specifically designed to avoid. The insistence on
processing fields on their natural geometric domain, rather than on a
flattened approximation, is precisely an insistence on admissibility
preservation: the representational geometry must be compatible with the
physical geometry.

\subsection{Entropy as Refinement Volume}

It is worth pausing on a potential confusion about the role of entropy in
\RSVP. The entropy density $S_t$ is not thermodynamic disorder in the
standard sense, nor is it Shannon information-theoretic entropy over
probability distributions. It is \textit{geometric admissibility volume}:
the local measure of how many fine-scale refinements are consistent with
the current coarse-scale state and accumulated history.

In a region of low entropy, the coarse-scale field $\varphi_t$ tightly
constrains the admissible fine-scale states: the compatibility conditions
at that location leave very little room. Reconstruction of fine-scale
detail from coarse inputs is nearly deterministic in such regions, because
almost all admissible refinements are mutually close. In a region of high
entropy, the coarse-scale structure leaves the fine-scale states largely
unconstrained: many very different fine-scale configurations are admissible,
and reconstruction becomes an act of selection among a large admissible
set rather than a nearly unique recovery.

This reinterpretation has a precise implication for the \FSA\ architecture.
The regions where the climate model has low reconstruction uncertainty ---
where its predictions are most reliable and consistent --- correspond to
low-entropy regions in the \RSVP\ sense: regions where the global constraint
structure tightly regulates the admissible fine-scale states. The regions
where the model has high uncertainty correspond to high-entropy regions:
where the constraint structure leaves many equally admissible fine-scale
configurations. The architecture's zero-shot super-resolution capability is
strongest precisely in the low-entropy regime, where constraint transport
alone is sufficient to determine the admissible refinements.

\subsection{Scale Hierarchy as Constraint Cascade}

The \RSVP\ framework treats scale decomposition as a constraint cascade
rather than a feature hierarchy. At the coarsest scale, the scalar field
$\varphi$ encodes global accessibility structure: the long-range coherence
that determines which large-scale trajectories are admissible. At
progressively finer scales, the vector field $\mathbf{v}$ and entropy
density $S$ encode local perturbative refinements of that structure. A
region of low entropy ($S$ small) is a region where coarse structure
tightly constrains fine structure: few refinements are admissible, and the
compression is nearly lossless. A region of high entropy ($S$ large) is a
region where coarse structure leaves fine structure relatively unconstrained:
many refinements are admissible, and the compression is inherently lossy.
This is not a limitation to be overcome; it is a structural feature of the
physical system.

\subsection{Constraint Geometry Is Not Reducible to Optimization}

Many contemporary theoretical frameworks describe physical and cognitive
systems as optimization processes: free-energy minimisation, variational
inference, least-action trajectories, or energy-landscape descent. The
\RSVP\ framework is compatible with such descriptions but is not reducible
to them, and the distinction is important for reading the empirical results
correctly.

Optimization presupposes a space over which optimization occurs. Constraint
geometry specifies the topology of that space itself: which trajectories
exist, which transitions are admissible, which refinements remain globally
consistent, and which regions of state space are structurally inaccessible
regardless of energetic cost. The admissibility function $\kappa(\omega,
H_t)$ is logically prior to any objective functional; it defines the
feasibility region over which objectives are then minimised.

A system may therefore occupy a locally low-energy configuration while
remaining globally inadmissible. Conversely, a trajectory may be
energetically expensive while remaining the unique admissible option given
the accumulated history. Constraint geometry regulates the structure within
which optimization takes place; it cannot itself be reduced to a particular
optimization objective without circularity.

This distinction becomes operationally visible wherever systems exhibit
long-horizon coherence that cannot be explained through local optimization
alone. The \FSA\ latent organization, the spin--valley Goldstone propagation,
and the trajectory-conditioning principles of \CLIO\ all depend on global
compatibility conditions that exceed local energetic minimisation. The
convergence of these systems toward geometry-aware design is not driven by
finding a better objective function; it is driven by the recognition that
the geometry of the feasibility space matters independently of any objective
defined over it.

\subsection{Degenerate Observables and Projection Equivalence}

A projection observable $\mathcal{O}: X \to Y$ is \textit{degenerate}
whenever dynamically distinct states in $X$ map to the same observable
state in $Y$:
\[
x_1 \neq x_2, \qquad \mathcal{O}(x_1) = \mathcal{O}(x_2).
\]
In such cases, observational equivalence does not imply dynamical
equivalence. Systems that appear identical on the observable surface may
possess fundamentally different admissibility structures and evolve
differently under perturbation.

Degeneracy is not a correctable measurement error; it is a structural
property of the projection. When the observable $\mathcal{O}$ eliminates
an entire subspace $X_{\mathrm{hid}}$ of dynamically significant degrees
of freedom, no improvement in the precision of $\mathcal{O}$-based
measurements can recover the lost structure. Recovery requires a
fundamentally different observational basis --- one that is not degenerate
with respect to $X_{\mathrm{hid}}$.

The spin--valley experiment demonstrates this directly. Charge observables
collapse dynamically distinct collective configurations into nearly
indistinguishable observable states; the hidden spin--valley modes are not
merely poorly measured but structurally eliminated by charge projection.
The resolution requires space-time imaging of the spin--valley fluctuations
directly, not a more precise charge measurement.

The same phenomenon appears in machine learning latent spaces. Two models
may achieve nearly identical endpoint accuracy while possessing radically
different latent trajectory organization and radically different
generalization behaviour. Observable equivalence at the metric level does
not imply equivalence of learned constraint structure. Benchmark
evaluations that rely exclusively on endpoint metrics are therefore
degenerate observables of the underlying representational system: they
collapse structurally distinct models into the same observable class.

% =============================================================
\section{The \CLIO\ Framework: Constraint-Leveraged Inference}
\label{sec:clio}

\subsection{The Inference Principle}

The \CLIO\ (Constraint-Leveraged Inference and Optimization) framework
extends the ontological commitments of \RSVP\ into a computational account
of how admissibility-preserving projection systems perform inference.
The central claim is that effective inference does not operate by
exhaustively searching the space of possible outputs and selecting the
most probable. It operates by first identifying the constraint structure
of the target domain and then projecting from that constraint structure
onto the space of admissible outputs.

Formally, the \CLIO\ inference pipeline is
\[
\omega \;\xrightarrow{\;\Pi\;} p \;\overset{\mathcal{R}}{\rightsquigarrow}\;
\hat{\ell} \;\xrightarrow{\;\mathcal{V}\;} \omega',
\]
where $\omega \in \Svalid$ is a constraint-valid structure, $\Pi$ is the
projection functor mapping it to a linguistic or computational
representation $p$, $\mathcal{R}$ is a stochastic rendering operator
producing an output $\hat{\ell}$, and $\mathcal{V}$ is a verification
map returning either a validated structure $\omega'$ or failure $\bot$.

The rendering operator $\mathcal{R}$ is stochastic rather than deterministic
because the projection $\Pi$ is in general many-to-one: multiple constraint
structures may project to the same representational point, and multiple
admissible outputs may be consistent with a single constraint structure.
The role of $\mathcal{V}$ is to check whether the rendered output lands
within the admissible region rather than on the boundary of the
representational manifold where the projection has become unreliable.

\subsection{Sparse Reconstruction and Latent Geometry}

The most distinctive \CLIO\ prediction is that effective inference systems
should exhibit sparse reconstruction: the ability to produce detailed,
coherent outputs from compressed constraint representations, without
exhaustive memorization of the fine-scale detail. This prediction follows
directly from the \RSVP\ account of entropy: in low-entropy regions where
constraint structure is tight, detailed reconstruction follows from the
constraints alone, without additional information. The constraint structure
is sufficient.

This is a strong claim, and one that the \FSA\ paper corroborates in a
specific domain. The zero-shot super-resolution capability of the \FSA\ is
precisely sparse reconstruction: the model produces fine-scale detail that
it has never directly observed by applying learned constraint compatibility
relations. The coarse state is sufficient to determine the admissible
fine-scale refinements; the model need not memorize those refinements
explicitly.

\begin{proposition}
\label{prop:sparse}
Let $\pi: X \to M$ be an admissibility-preserving projection, and let $x \in X$
be a state with $S(x) < \epsilon$ for some small $\epsilon > 0$. Then the
pre-image $\pi^{-1}(\pi(x))$ has volume bounded by a function of $\epsilon$
that tends to zero as $\epsilon \to 0$. In the limit of zero entropy,
the coarse representation determines the fine-scale state uniquely.
\end{proposition}

The proposition formalizes the intuition that low-entropy regions are
regions of tight constraint. The \FSA\ zero-shot super-resolution result
is an empirical demonstration that the trained architecture has learned
to identify these regions and exploit the tight constraints they impose.

\subsection{Latent Manifold as Trajectory Space}

The spontaneous organization of the \FSA\ latent space into cyclic seasonal
trajectories and historical warming drift is, from a \CLIO\ perspective, a
confirmation that the architecture has succeeded in preserving trajectory
geometry during compression. The latent manifold is not a cloud of
independently learned feature vectors; it is a geometric representation of
the dynamical system's phase space, complete with the qualitative structure
--- the cycles, the drift, the topological organization --- that
characterizes the underlying physics.

This is what \CLIO\ means by admissibility-preserving compression: the
compressed representation preserves not merely the instantaneous state but
the geometry of the space of admissible transitions between states. A
model that has learned this geometry can generalize to novel states and
novel time horizons in a way that a model organized around reconstruction
efficiency cannot.

\begin{definition}[Trajectory-Preserving Compression]
A compression $\pi: X \to M$ is \emph{trajectory-preserving} if, for every
admissible trajectory $\gamma: [0,T] \to X$ in the original space, the
projected trajectory $\pi \circ \gamma: [0,T] \to M$ is admissible in the
compressed space, and the mapping $\gamma \mapsto \pi \circ \gamma$ is
a homeomorphism from the space of admissible trajectories in $X$ to the
space of admissible trajectories in $M$.
\end{definition}

Trajectory-preserving compression is stronger than the basic
admissibility-preserving projection condition. The latter requires only
that individual admissible transitions survive compression; the former
requires that the global structure of the trajectory space be preserved.
The \FSA\ latent space organisation suggests the architecture achieves
something close to the stronger condition, at least within the training
distribution regime.

\subsection{The Inference Failure of Markov Systems}

The \CLIO\ framework provides a precise diagnosis of why standard
autoregressive systems fail at long-horizon coherence. A standard
next-step predictor operates without explicit constraint structure:
\[
P(x_{n+1} \mid x_n),
\quad H_t = \varnothing,
\quad \Pp_t = \varnothing,
\quad \mathcal{V} = \varnothing.
\]
Without history coupling, provenance continuity, or verification, the
output distribution at each step samples from the prior manifold of the
training distribution rather than from the admissible future of the current
trajectory. Short-horizon predictions may be locally accurate, but errors
accumulate because the system has no mechanism for maintaining global
constraint coherence.

The \FSA\ architecture avoids this failure by operating on temporally
structured compressed states rather than isolated frames. The coarse
manifold provides the global constraint that prevents the system from
drifting to locally plausible but globally inadmissible states. This is a
structural advantage over autoregressive approaches, and it reflects a
\CLIO\ inference principle rather than merely an engineering choice.

% =============================================================
\section{Formal Correspondence: \FSA\ as \RSVP/\CLIO\ Instance}
\label{sec:correspondence}

\subsection{The Mapping}

Having developed both the \FSA\ architecture and the \RSVP/\CLIO\
frameworks, we can now state the formal correspondence systematically.
The mapping is not exact --- the \FSA\ is an engineering system developed
without reference to the theoretical frameworks, and it lacks several of
their structural features --- but it is close enough to be theoretically
significant.

The coarse climate manifold $X_0$ corresponds to the low-frequency
component of the \RSVP\ scalar field $\varphi$: the global accessibility
structure that constrains all admissible fine-scale refinements. The
residual layers $R_k$ correspond to the perturbative refinement structure
of $\mathbf{v}$ and $S$: localized corrections that are constrained by
the coarser levels. The field-space operators $\mathcal{F}_{k \to j}$
correspond to constraint-preserving projection morphisms in the \CLIO\
pipeline. The compressed latent manifold corresponds to the entropy-weighted
accessibility geometry:
\[
S_{\mathrm{lat}} = \log \bigl|\A(X)\bigr|,
\]
where $|\A(X)|$ denotes the volume of the accessibility slice.

More precisely:

\begin{center}
\begin{tabular}{@{}>{\raggedright}p{5.0cm}>{\raggedright\arraybackslash}p{6.0cm}@{}}
\toprule
\textbf{Field-Space Autoencoder} & \textbf{\RSVP\ /\ \CLIO} \\
\midrule
Coarse climate manifold $X_0$ & Low-frequency accessibility field $\varphi$\\[0.3em]
Residual scale $R_k$ & Perturbative refinement of $\mathbf{v}$, $S$\\[0.3em]
Cross-scale operator $\mathcal{F}_{k\to j}$ & Approximation of projection functor $\Pi$ and verification $\mathcal{V}$\\[0.3em]
Latent manifold geometry & Entropy-weighted accessibility slice\\[0.3em]
Zero-shot super-resolution & Sparse reconstruction from constraint\\[0.3em]
Seasonal cycle in latent space & Admissible trajectory curvature\\[0.3em]
Historical warming drift & Irreversible shrinkage $\A_{t+1}^{\mathrm{con}}\subsetneq\A_t^{\mathrm{ctf}}$\\[0.3em]
Sphere-preserving processing & Admissibility-preserving projection\\
\bottomrule
\end{tabular}
\end{center}

\subsection{Operational Versus Structural Preservation}

The significance of the correspondence in the table above lies not in
identity of formalism but in convergence of architectural pressure. The
\FSA\ was not designed to instantiate admissibility geometry explicitly;
rather, admissibility-preserving structure emerged as a practical
requirement for stable multi-scale representation on physically constrained
domains. The table should therefore be read as mapping the engineering
solutions that pressure forced onto the theoretical concepts that the
\RSVP/\CLIO\ framework formalizes --- not as evidence that the engineers
were implementing the framework.

This distinction is the key to reading the correspondence correctly. An
architecture that \textit{operationally} preserves admissibility is one
whose outputs satisfy physical admissibility conditions because the training
data is organized around such conditions. The preservation is a consequence
of the training distribution, not a designed feature. An architecture that
\textit{structurally} preserves admissibility is one that explicitly encodes
the admissibility conditions and enforces them regardless of distributional
proximity to the training data. The \FSA\ achieves operational preservation;
the \RSVP/\CLIO/\TARTAN\ framework specifies what structural preservation
would require.

The practical consequence of this distinction becomes visible at distribution
boundaries. An operationally admissibility-preserving system generalizes
within the training distribution and begins to degrade at its edges.
A structurally admissibility-preserving system generalizes to any input
satisfying the same formal admissibility conditions as the training data,
regardless of its distributional proximity. When an operational system
encounters an out-of-distribution perturbation severe enough to push it
off the training manifold, it does not merely degrade gracefully; it
undergoes \textit{topological delamination} --- the compressed
representational trajectory separates from the physical admissibility
manifold entirely, because the system has no closed-loop mechanism for
detecting when it has left the admissible region. Structural preservation
would require precisely such a mechanism: a cohomological constraint
checker capable of computing the obstruction class $[\delta\sigma]$ and
flagging non-trivial values before they propagate through the model.

\subsection{The Projection Failure Theorem}

The \FSA\ paper's diagnosis of why flat convolutional architectures fail
on spherical domains corresponds to a general theorem about projection
failures in \RSVP.

\begin{theorem}[Projection Failure]
Let $\pi: S^2 \to \mathbb{R}^2$ be any continuous surjection from the
sphere to the plane. Then $\pi$ cannot be admissibility-preserving for
any physical field dynamics defined by rotational symmetry on $S^2$.
\end{theorem}

\begin{proof}[Sketch]
The failure is both topological and metric. Topologically, $S^2$ is compact
and boundaryless, while $\mathbb{R}^2$ is non-compact; by invariance of
domain, no continuous surjection $\pi: S^2 \to \mathbb{R}^2$ can be a
homeomorphism, so $\pi$ must either collapse distinct points or destroy the
boundary-free structure of $S^2$. Metrically, there exists no isometry
between $S^2$ (constant positive curvature) and $\mathbb{R}^2$ (zero
curvature), and no conformal equivalence either, since a conformal map would
require the Riemann sphere $\hat{\mathbb{C}} \cong S^2$ to be conformally
equivalent to $\mathbb{C}$, which it is not (they have different conformal
types). Consequently, any $\pi$ introduces unavoidable metric distortion.
The pushforward of the invariant vector fields generating $\mathrm{SO}(3)$
rotations on $S^2$ cannot map to a closed, smoothly invariant symmetry
group under the Euclidean transformation metric without introducing
coordinate singularities or metric divergence at at least one point. Since
rotationally symmetric dynamics on $S^2$ are defined by these invariant
vector fields, and their pushforward is metrically ill-behaved in
$\mathbb{R}^2$, any physically admissible transition along a great circle
geodesic maps to a transition that is generally inadmissible under the
Euclidean symmetry group. Admissibility cannot be preserved.
\end{proof}

This theorem is not a new result in differential geometry, but its
interpretation in the \RSVP\ framework gives it a new significance: flat
projection failure is not a numerical artifact to be managed by
preprocessing but a structural incompatibility between the representational
and physical geometries. The \FSA\ paper's insistence on processing on the
native geometric domain is therefore not an engineering preference but a
structural necessity.

\subsection{Latent Organization as Fixed-Point Signature}

The spontaneous organization of the \FSA\ latent space is, in \RSVP\
terms, the signature of a fixed-point structure in the compressed
accessibility geometry. A compression $\pi$ that preserves trajectory
geometry produces a latent manifold that is dynamically equivalent to the
original state space: the attractors, cycles, and drift directions of the
original dynamics appear as attractors, cycles, and drift directions in
the latent space.

The \RSVP\ fixed-point formalism (developed in the KES extension) treats
this as a stability condition:
\[
\KES^* = \mathcal{G}(\KES^*),
\quad \rho(D\mathcal{G}_{\KES^*}) < 1,
\]
where $\mathcal{G}$ is the recursive evolution operator and $\rho$ denotes
spectral radius. A compression is trajectory-preserving in the sense of
Definition~4.3 if and only if the compressed dynamics have a fixed-point
structure that is topologically equivalent to the original. The \FSA\
latent space organization is empirical evidence that the training process
has found such a compression.

\subsection{Where the Correspondence Breaks Down}

It is important to be precise about the limits of this correspondence.
The \FSA\ architecture achieves trajectory-preserving compression
operationally, in the sense that its compressed representations exhibit
the qualitative geometric properties expected of trajectory-preserving
systems. But it does not \textit{enforce} these properties explicitly. The
correspondence is a consequence of the architecture's geometric
commitments, not a designed feature.

More significantly, the \FSA\ lacks an explicit account of admissibility.
It learns which transitions are physically plausible from data, but it has
no formal structure for distinguishing between transitions that are
inadmissible due to physical law and transitions that are merely improbable
due to the training distribution. In \RSVP\ terms, the architecture
conflates $\Adm(X_t)$ with the high-probability region of $P(x_{t+1} \mid
x_t)$. For in-distribution inputs, these sets may be nearly identical.
For out-of-distribution inputs, they can diverge substantially, and the
architecture provides no mechanism for detecting or correcting this
divergence.

This is the primary theoretical gap that the \RSVP/\CLIO\ framework
addresses but the \FSA\ does not.

% =============================================================
\section{The \TARTAN\ Framework: Architectural Generalization}
\label{sec:tartan}

\subsection{From Instance to Architecture}

The \RSVP/\CLIO\ account of the \FSA\ explains why the architecture works
as well as it does and identifies the structural features that account for
its success. The \TARTAN\ (Trajectory-Aware Recursive Tiling with
Annotated Noise) framework generalizes these features into a principled
design vocabulary for admissibility-preserving architectures across domains.

\TARTAN\ takes the following as foundational design principles. Hierarchical
structure must be explicit: the architecture must specify the scale
decomposition rather than leaving it to be inferred. Residuals must be
semantically annotated: each residual correction must carry information
about its admissibility conditions, not merely its magnitude. Coarse-to-fine
reconstruction must be constraint-regulated: fine-scale detail must be
generated as what is admissible given coarse structure, not as what is
most probable given the training distribution. And trajectory conditioning
must be global: the architecture must maintain consistency with the full
accessible history, not merely with the immediately preceding state.

\subsection{Tiling and Residual Annotation}

The \TARTAN\ framework's central architectural commitment is to recursive
tiling with annotated noise. A coarse field $X_0$ tiles the domain into
regions of approximately homogeneous admissibility structure. Each tile is
then recursively subdivided, and each subdivision is assigned a residual
correction $R_k$ that carries semantic annotation: information about what
class of physical process the residual represents and what admissibility
constraints govern it.

Formally, the tile decomposition is
\[
\mathcal{D}_k: \Omg_t \to \bigl\{(U_i^{(k)}, R_i^{(k)})\bigr\}_{i},
\]
where $U_i^{(k)}$ is the $i$-th tile at scale $k$ and $R_i^{(k)}$ is its
annotated residual. The annotation carries both the semantic label of the
residual (which physical process it represents) and the admissibility
condition it must satisfy given the coarser tiles.

The \FSA\ architecture implements an unannotated version of this structure.
Its multi-scale residual decomposition tiles the domain into scale strata
and assigns cross-scale operators to the boundaries between strata. The
absence of explicit semantic annotation is one of the primary differences
between the \FSA\ implementation and the full \TARTAN\ design.

\subsection{Trajectory Conditioning and the Annotation Function}

The trajectory-conditioning component of \TARTAN\ is the feature that most
directly addresses the gap identified in the previous section. Rather than
conditioning each residual correction only on the current coarse state,
\TARTAN\ conditions each residual on the full admissible history:
\[
R_i^{(k)}(x, t) = f\bigl(X_0(t), H_t, \Psi_t, \kappa(\cdot, H_t)\bigr),
\]
where $\Psi_t$ is the constraint field and $\kappa(\cdot, H_t)$ is the
admissibility function. This ensures that fine-scale reconstructions are
not merely locally consistent with the coarse state but globally consistent
with the entire accessible history.

The \FSA\ implements a version of temporal conditioning through its
diffusion architecture, which operates over temporally structured compressed
states. This is structurally similar to \TARTAN\ trajectory conditioning,
but without the explicit admissibility constraint. The \FSA\ conditions on
the recent past; \TARTAN\ conditions on the admissible past, which is a
strictly stronger condition.

\subsection{Comparison with the Field-Space Operators}

The \FSA\ field-space operators $\mathcal{F}_{k \to j}$ are the
architectural component that most closely resembles the \TARTAN\ residual
annotation machinery. In both cases, cross-scale operators take multiple
scale strata as input, ensuring that cross-scale compatibility is a
multi-way rather than pairwise constraint. In both cases, the operators
are learned rather than specified, but the architectural context constrains
their function.

The primary difference is that the \FSA\ operators are learned from
reconstruction objectives, while the \TARTAN\ annotation function is
trained to satisfy explicit admissibility criteria. The \FSA\ operators
will converge to whatever cross-scale compatibility structure minimizes
reconstruction loss on the training distribution. The \TARTAN\ annotation
function is explicitly constrained to respect the admissibility conditions
derived from the physical or structural domain. In practice, on
in-distribution data, these may converge to similar solutions. On
out-of-distribution data, they will diverge in ways that reflect the
presence or absence of explicit admissibility enforcement.

\subsection{Why Constraint Hierarchies Generalize}

The zero-shot super-resolution capability of the \FSA\ is an instance of a
more general principle: constraint hierarchies generalize across resolution
and domain in ways that direct input-output mappings do not. The reason is
structural. A direct mapping from coarse input to fine output encodes the
statistical regularities of the training distribution at the specific
resolution pair presented during training. It has no mechanism for
extrapolating to resolution pairs it has not seen. A constraint hierarchy,
by contrast, encodes the compatibility conditions between scale strata: what
fine-scale structures are admissible given what coarse-scale states, and
why. These compatibility conditions are properties of the physical system,
not of the training distribution. They hold at resolution strata that were
never presented during training, because the physics does not change when
the measurement resolution does.

This is the \TARTAN\ generalization of the \FSA\ zero-shot result. Any
architecture that learns constraint compatibility rather than direct
input-output mappings should exhibit zero-shot generalization to novel
constraint configurations. The prediction is domain-independent: it applies
equally to protein folding dynamics (where the constraints are
thermodynamic admissibility conditions on backbone geometry), to turbulent
flow (where the constraints are Navier--Stokes admissibility conditions
across Reynolds number regimes), and to semantic coherence in language
(where the constraints are logical and narrative admissibility conditions
that hold across surface paraphrase variants). In each case, the
generalizing object is the constraint structure, not the surface form.

\subsection{Generalizing Zero-Shot Capability}

One of the most practically significant features of the \FSA\ --- its
zero-shot super-resolution capability --- follows naturally from the
\TARTAN\ design principles. The capability arises because the architecture
has learned the constraint compatibility relations between scale strata
independently of the specific resolution targets used during training.
Given those relations, it can apply them to novel resolution combinations.

The \TARTAN\ framework predicts this capability as a general consequence
of design that preserves admissibility. Any architecture that learns
constraint compatibility rather than direct input-output mappings should
exhibit zero-shot generalisation to novel constraint configurations. This
is a testable prediction extending beyond the climate domain. An architecture
trained on constraint-compatible multi-scale representations of protein
folding dynamics, or of turbulent flow, or of language at multiple
granularities, should exhibit analogous zero-shot capability in its
respective domain.

% =============================================================
\section{Sheaf-Theoretic Formalization and the Mathematical Horizon}
\label{sec:sheaf}

\subsection{From Engineering Constraints to Cohomology}

The \FSA\ field-space operators enforce cross-scale compatibility
operationally: the architecture is designed so that outputs at finer
resolutions are consistent with inputs at coarser resolutions, by
construction. This engineering constraint has a precise mathematical
expression. Whenever a system must enforce local conditions that are
simultaneously consistent with all other local conditions across a covered
domain, the natural language for that requirement is sheaf theory.

A sheaf is exactly the data structure that encodes local-to-global
compatibility: it assigns information to each open set in a cover, and
requires that the assignments on overlapping sets agree on their
intersection. The \FSA\ cross-scale operators are engineering implementations
of sheaf sections: they assign field representations to each scale stratum,
and the multi-strata input requirement ensures that these assignments are
mutually consistent. The shift to sheaf language is therefore not an
imported mathematical metaphor but an explicit formalization of the same
compatibility structure already implemented operationally by the field-space
operators. The categorical machinery makes visible the structure that the
engineering is already enforcing.

This also explains why the zero-shot super-resolution capability has a
natural cohomological interpretation. The capability corresponds to the
ability to extend a cocycle from observed scale strata to unobserved ones.
That extension is trivial --- uniquely determined --- precisely when the
local compatibility conditions at the observed scales are tight enough to
constrain the admissible extensions. The sheaf language turns a benchmark
result into a structural theorem about the cohomological complexity of the
domain.

\subsection{Local-Global Consistency as Sheaf Condition}

The compatibility constraints that govern the \TARTAN\ tiling structure ---
and that the \FSA\ implements operationally --- can be given a precise
mathematical formulation in terms of sheaf theory. A sheaf over the scale
hierarchy encodes the requirement that local admissibility conditions be
globally consistent: what is admissible at a fine scale must be compatible
with what is admissible at all coarser scales, and the compatibility
conditions must satisfy a gluing axiom.

Let $\mathfrak{U} = \{U_i\}$ be an open cover of the domain $\fM$ by
admissibility regions. The local accessibility sheaf $\scrA$ assigns to
each open set $U_i$ the admissible set of transitions within that region:
\[
\scrA(U_i) = \Adm(X_t) \cap U_i.
\]
The gluing condition requires that locally admissible transitions be
globally consistent:
\[
\scrA(U_i)\big|_{U_i \cap U_j} = \scrA(U_j)\big|_{U_i \cap U_j}.
\]
When this condition fails, the system has a \emph{global section
obstruction}: there is no globally admissible trajectory consistent with
the local admissibility conditions at each scale.

\subsection{Obstruction Classes and Irreducible Inconsistency}

The obstruction to finding a global section is measured by the first
\v{C}ech cohomology class:
\[
[\delta\sigma] \in \check{H}^1(\mathfrak{U}, \scrA).
\]
Vanishing cohomology $[\delta\sigma] = 0$ corresponds to a globally
consistent history: the local admissibility conditions are compatible and
a global section exists. Non-trivial cohomology $[\delta\sigma] \neq 0$
encodes an irreducible inconsistency: the local conditions are mutually
incompatible in a way that cannot be resolved by refining the cover.

In the climate domain, this obstruction structure has a physical
interpretation. A non-trivial cohomology class would correspond to a
climate state in which the locally admissible refinements at different
scales are globally inconsistent --- a situation that the physical system
cannot actually reach. The \FSA\ architecture avoids such states by
construction, because its field-space operators enforce cross-scale
compatibility. But it does so without explicitly computing or checking the
cohomology class.

The \RSVP/\TARTAN\ framework provides a principled path toward explicit
obstruction detection: an architecture equipped with a cohomological
constraint checker could identify globally inadmissible states before
generating them, rather than relying on the training distribution to
have excluded them.

\subsection{Holonomy and Path Dependence}

A deeper sheaf-theoretic feature of the \RSVP\ framework is the holonomy
of the constraint connection. If the constraint field $\Psi_t$ is treated
as a connection on the admissibility bundle over $\fM$, then the holonomy
of a closed loop $\gamma$ in the model space is
\[
\mathrm{Hol}_\gamma(\Psi) = \mathcal{P}\exp\!\Bigl(
\oint_\gamma A_\Psi \, d\ell \Bigr) \in \mathrm{Aut}(\Omg_t),
\]
where $A_\Psi$ is the connection form associated with $\Psi$. Non-trivial
holonomy implies that the admissibility structure is path-dependent: the
set of admissible transitions depends not only on the current state but on
the history of states through which the system has passed.

This is the formal expression of the intuition, central to \RSVP, that
history matters intrinsically rather than merely instrumentally. The \FSA\
architecture implicitly exhibits path dependence through its temporal
conditioning --- the compressed state at time $t$ depends on the sequence
of compressed states that preceded it, not merely on the last state. But
it lacks an explicit account of the holonomy structure that generates this
dependence.

\subsection{Refinement Cohomology}

The zero-shot super-resolution capability of the \FSA\ suggests a
cohomological interpretation of the reconstruction process. The
admissibility conditions at different scales define a cochain complex:
\[
\cdots \to C^{k-1}(\mathfrak{U}, \scrA) \xrightarrow{\delta}
C^k(\mathfrak{U}, \scrA) \xrightarrow{\delta} C^{k+1}(\mathfrak{U}, \scrA) \to \cdots
\]
A super-resolution query at a novel resolution $k^*$ corresponds to
extending a cocycle from the scales $\{0, 1, \ldots, k^*-1\}$ to a cocycle
that includes $k^*$. The query is answerable if and only if the extension
is cohomologically trivial: if the admissibility conditions at scales
$0$ through $k^*-1$ uniquely determine an admissible refinement at scale
$k^*$.

The \FSA\ zero-shot super-resolution result is empirical evidence that the
training process has learned a representation in which this extension is
approximately trivial within the training distribution. The \TARTAN\
framework predicts that explicitly enforcing the cohomological constraint
during training would extend the zero-shot capability beyond the training
distribution.

\subsection{Constraint Frustration and Obstruction Structure}

Non-trivial cohomology corresponds physically to \textit{frustrated}
constraint systems: collections of locally admissible conditions that cannot
be simultaneously satisfied globally. Spin glasses, geometrically frustrated
magnets, and incompatible boundary conditions in fluid dynamics all exhibit
this structure. In each case, every local subsystem has a well-defined
preferred state, but the local preferences are mutually incompatible across
the full domain.

The significance of the sheaf-theoretic formalism is that it treats such
failures not as numerical instability or optimization failure but as
\textit{topological obstruction}. A non-vanishing obstruction class
$[\delta\sigma] \neq 0$ does not indicate that the optimizer has not
converged or that more training data is required. It indicates that no
globally admissible section exists for the local conditions supplied:
the incompatibility is structural and cannot be resolved by any continuous
local adjustment.

This provides a principled distinction between two failure modes that
existing machine learning systems frequently conflate. An optimization
failure is recoverable: with more capacity, more data, or a better
objective, the system can approach the globally consistent solution. A
topological obstruction is irrecoverable within the given constraint
structure: the system must change its admissibility conditions (i.e., the
physical model) rather than its optimization. Architectures lacking an
explicit obstruction formalism cannot distinguish between these cases.
They interpret topological failures as optimization failures, add capacity,
and discover that the capacity has no effect --- because the problem is not
insufficient optimization but structural inadmissibility of the local
conditions.

% =============================================================
\section{Spin--Valley Modes in Moiré Superlattices: Hidden Constraint Propagation}
\label{sec:moire}

\subsection{The Physical Setting}

Moiré superlattices formed by stacking two atomically thin layers of
transition metal dichalcogenides at a small twist angle constitute one of
the richest experimental platforms in contemporary condensed matter physics.
The periodic interference pattern between the two lattices creates an
effective superlattice potential with a unit cell orders of magnitude larger
than the atomic scale, and the resulting flat bands concentrate kinetic
energy suppression in precisely the regime where interaction effects become
dominant. Twisted WSe$_2$ bilayers, in particular, have recently been shown
to host superconductivity~\cite{Xia2025,Guo2025}, correlated insulator
phases, and topological states, making them a setting where multiple
competing orders --- excitonic, magnetic, superconducting --- are
simultaneously accessible.

The paper we treat as our second empirical reference~\cite{Xiong2026}
studies a distinct but related phenomenon: the propagation of collective
spin--valley modes in a moiré-hosted spin--valley superfluid. The
theoretical background was established by Bi and Fu~\cite{BiFu2021}, who
predicted an excitonic density wave and spin--valley superfluid in bilayer
transition metal dichalcogenides arising from interlayer Coulomb attraction
between electrons in opposite spin--valley sectors. The key feature of this
state is that it supports a charge-neutral order parameter: the condensed
excitons carry no net charge, only correlated spin and valley quantum
numbers.

\subsection{Hidden Does Not Mean Unphysical}

The language of hidden fields and hidden manifolds risks a misreading that
must be addressed directly. Throughout this essay, ``hidden'' does not mean
metaphysically inaccessible, unobservable in principle, or existing in a
realm beyond the physical. The spin--valley modes are hidden only relative
to a particular measurement basis: ordinary charge-sensitive probes.
Once the measurement strategy changes --- from static charge response to
direct space-time imaging of spin--valley fluctuations --- the previously
hidden structure becomes fully and precisely empirically accessible. The
Xiong et al. paper does exactly this: it makes the hidden field visible by
choosing a non-degenerate observational basis.

The \RSVP\ framework makes the same ontological claim. It does not posit
an inaccessible substrate beyond observation. It argues that different
observational projections preserve different subsets of the underlying
admissibility structure. A projection may be operationally convenient ---
charge measurement is simpler than space-time spin--valley imaging --- while
still eliminating the dynamically relevant organizational degrees of freedom.
The hidden field is hidden relative to a representational basis that was
optimised for different observables, not hidden in any absolute or
metaphysical sense.

This matters for interpreting the broader convergence thesis. The claim
that ``hidden dynamical manifolds are more physically fundamental than
observable state variables'' is not a mystical claim about hidden realities.
It is a structural claim about which representational choices preserve the
information necessary for long-horizon prediction and global constraint
coherence. The manifest structure of any projection is real; the claim is
that it is not \textit{sufficient} for the physical purposes under discussion.

\subsection{Observable Surfaces and Hidden Dynamics}

The experimental and theoretical situation in twisted WSe$_2$ presents with
unusual clarity a distinction that the \RSVP\ framework treats as
structurally fundamental: the distinction between the observable projection
surface and the hidden constraint dynamics beneath it. Standard condensed
matter measurement tools --- transport, optics, scanning probes ---  are
sensitive to charge. They record the charge density, the charge current, the
charge compressibility. These are real physical quantities, and measuring
them carefully yields genuine information about the system.

The difficulty is that the most important collective organization of the
spin--valley superfluid state is charge-neutral. The order parameter carries
no net charge; the Goldstone modes that propagate through the system carry
no charge current; the Higgs fluctuations that stabilize the condensate
produce no charge response. The observable surface is real, but it is not
the physical system. The physical system includes both the observable
projection and the hidden constraint geometry that the projection does not
access.

This situation is formally identical to the \RSVP\ decomposition
$X_t = X_t^{\mathrm{obs}} \oplus X_t^{\mathrm{hid}}$ with
$\mathcal{O}(X_t^{\mathrm{hid}}) = 0$. The observable projection $\mathcal{O}$
kills precisely the degrees of freedom that carry the system's collective
organization. Recovering those degrees of freedom requires a measurement
strategy that is sensitive to the hidden field --- one that does not project
onto the observable surface first and attempt reconstruction from that
projection, but that accesses the constraint geometry directly.

\subsection{Collective Modes as Constraint Propagation}

The Goldstone modes of the spin--valley superfluid are not merely soft
fluctuations associated with a broken symmetry. They are the physical
realization of constraint propagation: the mechanism by which the local
admissibility conditions imposed by the moiré superlattice potential
propagate coherently across the sample and maintain long-range collective
order.

In the language of \RSVP, the moiré superlattice potential plays the role
of the constraint field $\Psi_t$: it specifies, at each spatial location,
which configurations of the spin--valley degrees of freedom are energetically
admissible. The spin--valley superfluid state is the global section of the
admissibility sheaf $\scrA$ over the moiré domain: the unique (up to gauge)
configuration that is locally admissible everywhere simultaneously. The
Goldstone modes are the gapless perturbations of this global section that
cost no energy to the leading order, because they correspond to global
rotations of the order parameter that respect all the local admissibility
conditions uniformly. The Higgs mode is the perturbation that costs energy
because it locally violates the amplitude constraint: it attempts to change
$|\psi|$ from its condensate value, which the local admissibility condition
resists.

The propagation velocity of the Goldstone mode is therefore not merely a
material parameter but a measure of the stiffness of the global constraint
structure: how strongly the local admissibility conditions resist relative
phase fluctuations between spatially separated regions of the sample. A
high propagation velocity indicates a stiff constraint structure, one where
the local admissibility conditions are tightly coupled across space. A low
velocity indicates a soft constraint structure, where local admissibility
conditions are relatively independent.

\subsection{The Measurement Problem and Its Resolution}

The central experimental challenge is that the spin--valley modes are
charge-decoupled. Ordinary transport measurements, which detect charge
currents, are insensitive to the propagating excitations because the condensed
excitons carry no net charge. Optical measurements that probe charge density
are similarly blind. Standard steady-state measurements show only the total
response, in which charge-coupled excitations dominate and the
charge-neutral spin--valley modes are eclipsed entirely.

The experimental resolution is a pump--probe space-time imaging technique.
A shaped pump laser injects a localized packet of excitations at time $t=0$;
a delayed wide-field probe then tracks their spatial evolution as a function
of time delay $\Delta t$, recording the full space-time trajectory rather
than a static integrated response. Near the van Hove singularity --- a region
of the moiré phase diagram with particularly high density of electronic states
--- the technique resolves two propagating collective modes. The modes carry
\textit{opposite} spin--valley currents and propagate at drastically different
velocities: one faster than $3\;\mathrm{km\,s^{-1}}$ with partially
ballistic behaviour consistent with a gapless excitation, the other slower
and diffusive, consistent with a gapped excitation. The fast, ballistic mode
is interpreted as the Goldstone mode of the intervalley coherent (IVC)
ground state; the slow, diffusive mode is interpreted as the amplitude
(Higgs) mode with an energy gap. The authors describe their mission as
developing the capacity to directly record the ``complete space-time
evolution'' of many-body quantum systems --- language that is
almost exactly the \RSVP\ prescription for trajectory recovery rather than
state-snapshot measurement.

\subsection{\RSVP\ Interpretation of Phase and Amplitude Modes}

The Goldstone--Higgs mode decomposition in the spin--valley superfluid maps
naturally onto the \RSVP\ field decomposition $X_t = (\varphi_t,
\mathbf{v}_t, S_t)$. The phase (Goldstone) modes correspond to the
vector flow $\mathbf{v}_t$: they govern coherent directional propagation
through the constraint manifold, are massless (gapless), and carry
long-range organizational information about the system's collective state.
The amplitude (Higgs-like) modes correspond to fluctuations of the scalar
density $\varphi_t$: they are energetically costly, localized in effect,
and reflect the local cost of deviating from the condensate density.

More formally, let $\psi = \sqrt{\rho}\, e^{i\theta}$ denote the
spin--valley order parameter in polar decomposition, where $\rho$ is the
condensate density and $\theta$ is the phase. The effective action in the
vicinity of the condensate has the form
\[
\mathcal{L} \approx
\tfrac{1}{2}(\partial_\mu \theta)^2
+ \tfrac{1}{2}m_H^2 (\delta\rho)^2
+ \cdots,
\]
where the first term governs gapless phase propagation and the second
governs gapped amplitude fluctuations with Higgs mass $m_H \propto
\sqrt{\rho_0}$. In \RSVP\ notation:
\begin{align}
\text{Phase modes:} &\quad
\mathbf{v}_t \sim \nabla\theta, \quad
\partial_t\theta = -\nabla\varphi + \mathcal{T}(\mathbf{v}_t),\\
\text{Amplitude modes:} &\quad
\delta\varphi_t \sim \delta\rho, \quad
\partial_t(\delta\varphi) = -m_H^2 \delta\varphi + D_\varphi\nabla^2(\delta\varphi).
\end{align}
The phase modes preserve coherence and carry the system's organizational
information across space; the amplitude modes stabilize the local condensate
density against perturbative fluctuations. This is precisely the \RSVP\
distinction between low-entropy coherence-preserving flows (which
govern accessibility geometry) and local energetic stabilization
perturbations (which regulate the scalar field against collapse).

\subsection{Hidden Modes and Observable Projections}

The most theoretically significant feature of the spin--valley experiment is
the explicit demonstration that the dynamically relevant structure of the
system is not accessible through its standard observable projection. This is
the condensed matter version of the \RSVP\ critique of purely observational
epistemology, and it is worth stating precisely.

Let $\mathcal{O}$ denote the observable projection: the map from the
physical state $X_t$ to the set of measurable quantities (charge density,
conductance, optical response). The charge-decoupled spin--valley modes lie
in the kernel of $\mathcal{O}$: they contribute nothing to the observable
signal. The physically significant dynamics --- the propagating Goldstone
modes that organize the collective state --- are, in this sense, invisible.
They belong to the hidden part of the field,
\[
X_t = X_t^{\mathrm{obs}} \oplus X_t^{\mathrm{hid}},
\quad \mathcal{O}(X_t^{\mathrm{hid}}) = 0.
\]

This decomposition is not merely a technical accident of the measurement
apparatus. It reflects a structural feature of the physics: the symmetry
that protects the spin--valley modes from charge hybridization is the same
symmetry that makes them dynamically significant. The hidden modes are
hidden precisely because they are charge-neutral; and they are the carriers
of collective order precisely for the same reason.

In \RSVP\ terms, the observable projection $\mathcal{O}$ plays the role of
a lossy compression that eliminates admissibility-relevant structure from
the representation. The standard condensed matter toolkit --- transport,
optics --- operates exclusively on $X_t^{\mathrm{obs}}$ and is therefore
blind to the constraint dynamics of $X_t^{\mathrm{hid}}$. Recovering those
dynamics requires a measurement strategy that is sensitive to the hidden
field: direct space-time imaging of the spin--valley fluctuations rather than
inference from charge observables.

This is the physical analogue of the \CLIO\ prescription for
constraint-leveraged inference: do not project the constraint structure
onto a lower-dimensional observable manifold and attempt reconstruction
from that projection. Instead, develop measurement or inference strategies
that are sensitive to the admissibility structure directly.

\subsection{Space-Time Imaging as Trajectory Recovery}

The experimental technique developed in the moiré paper is a form of
trajectory recovery. Rather than recording a static snapshot of the
collective state, the authors record its complete space-time evolution: how
the spin--valley fluctuations propagate across the sample as a function of
both position and time. The result is a direct observation of the
trajectories of the collective modes rather than merely their instantaneous
configuration.

This maps onto the \CLIO\ argument about Chain-of-Memory versus static
symbolic traces with unusual precision. The authors are not asking ``what
is the charge density at time $t$?'' --- the static observable projection.
They are asking ``how does the spin--valley order parameter evolve from
$t_0$ to $t$?'' --- the trajectory in the hidden field space. The
physically meaningful answer is the trajectory, not the snapshot.

The \RSVP\ formalization of this distinction is the provenance structure
$\Pp_t = \Pp(\omega_t)$. Provenance records the history of admissible
transitions that produced the current state: it is the trajectory, not
the state, that carries the organizational information. A system whose
collective modes propagate coherently over macroscopic distances is a
system whose provenance trace extends over those distances; the coherence
is a property of the trajectory rather than of any instantaneous
configuration.

\subsection{Emergent Macroscopic Organization from Local Constraints}

A further feature of the spin--valley experiment resonates with the
\RSVP/\TARTAN\ framework: the emergence of macroscopic collective
organization from purely local interaction constraints. The moiré potential
provides a periodic local constraint --- the superlattice --- that
frustrates the electrons into a correlated ground state. The spin--valley
superfluid emerges not from any single interaction but from the global
consistency of local admissibility conditions across the entire sample.

This is the physical realization of the sheaf-theoretic structure developed
in Section~\ref{sec:sheaf}. The local admissibility condition at each moiré
unit cell is the constraint that interlayer excitons occupy the
energetically favoured spin--valley sectors. The global section is the
ordered condensate state. The Goldstone modes are the gapless fluctuations
of the global section; the Higgs mode is the cost of a section that
violates the local amplitude constraint.

The experimental observation that these modes propagate coherently across
the sample --- rather than decaying into local incoherence --- is evidence
that the sheaf gluing condition is satisfied: the local admissibility
conditions are globally consistent, and a global section exists. The
absence of a propagating Goldstone mode would correspond to a non-trivial
\v{C}ech cohomology class: the local constraints would be mutually
incompatible in a way that prevents the formation of long-range collective
order.

% =============================================================
% =============================================================
\section{Novel Theoretical Conclusions}
\label{sec:novel}

The formal correspondences developed in the preceding sections are not
merely analogical. When pressed, they yield conclusions that are more
specific than the general thesis of convergent rediscovery, and that have
implications extending beyond either empirical paper. This section states
five such conclusions explicitly.

\subsection{Compression and Measurement as Formal Duals}

The \FSA\ paper studies a compression problem: reduce a high-dimensional
physical field to a compact latent representation while preserving
admissibility structure. The moiré experiment studies a measurement
problem: reconstruct hidden collective dynamics from an observable
projection surface. These are dual operations in a precise sense.

Let $\pi: X \to M$ denote the compression map and
$\mathcal{I}: M_{\mathrm{obs}} \to X_{\mathrm{hid}}$ denote the
reconstruction map from observables to hidden dynamics. The two papers
demonstrate, from opposite directions, that both operations must respect
the same geometric constraint: the admissibility structure of the physical
system. The \FSA\ shows that compression fails when it violates the
geometry of the underlying process. The moiré experiment shows that
measurement fails when it projects onto a surface that eliminates the
physically significant degrees of freedom.

The formal duality is:
\[
X \;\xrightarrow{\;\pi\;}\; M
\qquad\text{(compression: preserve admissibility downward)}
\]
\[
M_{\mathrm{obs}} \;\xrightarrow{\;\mathcal{I}\;}\; X_{\mathrm{hid}}
\qquad\text{(measurement: recover admissibility upward)}
\]
Both operations succeed when they respect constraint geometry and fail
when they do not. This suggests a unified principle: \textit{effective
measurement and effective compression obey the same admissibility-preserving
geometric constraints}. The design of good scientific instruments, the
design of good latent representations, and the design of good inference
systems may all reduce to the same underlying problem.

\begin{proposition}[Compression--Measurement Duality]
Let $X$ be a field space with admissibility structure $\Adm(X)$. A
compression $\pi: X \to M$ is admissibility-preserving iff the
reconstruction map $\mathcal{I} = \pi^{-1}|_{\Adm}$ exists and maps
$M$ back into $\Adm(X)$. The existence of $\mathcal{I}$ is equivalent
to $\pi$ being injective on $\Adm(X)$: the admissible set is not
collapsed by the compression.
\end{proposition}

The proposition states that a compression is good --- in the
admissibility-preserving sense --- if and only if it supports a good
reconstruction of the admissible region. The two operations are not merely
analogous; they are inverses, and their quality is determined by the same
geometric property of the map $\pi$.

\subsection{Hidden Dynamical Manifolds as Primary}

Both papers provide empirical evidence for a claim that is stronger than
the convergence thesis: that hidden dynamical manifolds may be more
physically fundamental than the observable state variables that standard
theory and measurement treat as primary.

The moiré experiment demonstrates this in the clearest possible form.
The physically significant organization of the spin--valley superfluid ---
the collective order, the propagating modes, the broken symmetry that
defines the phase --- exists entirely in the hidden field $X_{\mathrm{hid}}$.
The observable surface $X_{\mathrm{obs}}$ carries none of it. An account of
the system that begins with observable state variables and attempts to derive
the collective organization from them is not merely incomplete; it is
organized around the wrong primary objects.

The \FSA\ result makes the same point computationally. The Field-Space
architecture, trained only on residuals and never exposed to the coarse
seasonal and warming signal, spontaneously recovers the full dynamical
geometry of the climate system in its latent space. The convolution baseline,
trained on the full field including the large-scale climatological structure,
loses it. The organized latent geometry emerges from the constraint structure
of the residuals, not from the surface statistics of the observable field.

Together these results support the \RSVP\ ontological claim:
\[
X_{\mathrm{obs}} = \Pi(X_{\mathrm{constraint}}),
\]
where the observable field is a secondary projection of the hidden constraint
manifold rather than a primary object. This is structurally close to the
programme of generalized effective field theory, where macroscopic
observables are understood as integrations-out of microscopic degrees of
freedom. The difference is that in the \RSVP\ account, the hidden degrees
of freedom are not microscopic in the conventional sense but
\textit{constraint-geometric}: they are the admissibility conditions that
govern which trajectories are physically realizable, rather than high-energy
modes that happen to be integrated out.

\subsection{Deep Learning as Implicit Field Theory}

The \FSA\ architecture, examined from the perspective of what it actually
computes, is structurally much closer to discretized field evolution than to
classical symbolic AI or to the lookup-table intuition that motivates most
large-scale ML. It preserves topology by operating on the HEALPix sphere.
It evolves latent states continuously through temporally structured diffusion.
It transports compatibility constraints across scale boundaries rather than
mapping surface statistics. It encodes trajectory organization rather than
point-to-point correspondences.

These are precisely the structural commitments of a numerical field theory.
The architecture's success --- its superior compression efficiency, its
organized latent space, its zero-shot generalization --- depends directly
on respecting geometry, preserving continuity, and maintaining cross-scale
constraint consistency: the properties that define well-posed field theories.
The architecture succeeds to the extent that it behaves like a dynamical
field, and fails to the extent that it behaves like a lookup system.

A more precise formulation: under optimization pressure from physically
constrained domains, deep learning architectures increasingly become
discretized field theories without explicitly recognizing or naming this
transition. The \FSA\ is a data point in what appears to be a systematic
empirical trend. GraphCast~\cite{Lam2023}, Pangu-Weather~\cite{Bi2023},
FourCastNet~\cite{Kurth2023} --- the leading neural weather prediction
systems --- all converge on architectures that preserve spherical geometry,
propagate information continuously across spatial scales, and maintain
temporal coherence through structured latent dynamics. The transition from
storage-first to field-first representation is being driven by physics, not
by theoretical preference.

\subsection{Irreversibility as Geometric Contraction}

Both papers contain evidence that the physically important object is not
instantaneous state but constrained historical trajectory, and that
irreversibility is not primarily a thermodynamic property but a geometric one.

The \FSA\ latent space visualization is particularly revealing. Encoded
daily climate states from 1940 to 2024 drift smoothly and \textit{directionally}
through the latent manifold when coloured by year, with earlier states
concentrated in one region and more recent states displaced along an
outer arc. This drift is not a random walk; it has a consistent direction
in latent space that mirrors the directional character of global warming.
The trajectory through the latent manifold is irreversible not because
entropy increases in the thermodynamic sense, but because the accessible
region of the latent manifold at late times is a proper subset of the region
accessible at early times: certain latent states that were reachable in 1940
have become unreachable by 2024, and not merely because they are unlikely
but because the constraint structure of the system has shifted.

This is the \RSVP\ account of temporal asymmetry:
\[
\A_{t+1}^{\mathrm{consistent}} \subsetneq \A_t^{\mathrm{counterfactual}},
\]
stated in empirical rather than formal terms. The latent manifold contracts
under historical constraint accumulation. The contraction is geometric ---
a reduction in the volume of the accessible region --- rather than
thermodynamic. Irreversibility is not the growth of disorder but the
shrinkage of possibility.

The moiré experiment provides a physical analogue: the Goldstone modes
propagate directionally, carrying spin--valley current in a specific direction
across the sample. The propagation is not time-reversible in the
experimental sense: exciting a localized packet and watching it spread
produces a specific asymmetric space-time evolution. The constraint geometry
of the IVC state selects a direction of propagation; the trajectory through
the collective mode space is path-dependent in exactly the sense the \RSVP\
framework predicts.

\subsection{A Generalization Theorem}

The \FSA\ zero-shot capability across both resolution domains (HEALPix
levels 6--8) and model domains (ERA5 to MPI-ESM at 100~km resolution)
supports a general theoretical claim that is stronger than the observation
of a single benchmark result.

\begin{theorem}[Generalization from Constraint Geometry]
\label{thm:gen}
Let $\pi: X \to M$ be an admissibility-preserving compression and let
$\mathcal{C}$ denote the constraint compatibility structure learned during
training on distribution $\mathcal{D}$. Then the system generalizes
zero-shot to any input $x \notin \mathcal{D}$ satisfying the same
admissibility conditions as $\mathcal{D}$, in the sense that the
reconstruction $\pi^{-1}(\pi(x))$ lands in $\Adm(X)$.
\end{theorem}

\begin{proof}[Sketch]
If $\pi$ is admissibility-preserving, then $\pi$ maps $\Adm(X)$ injectively
into $M$. The learned compatibility structure $\mathcal{C}$ is a
representation of this map restricted to the training distribution
$\mathcal{D} \cap \Adm(X)$. For any $x \in \Adm(X)$ (regardless of
whether $x \in \mathcal{D}$), the admissibility conditions satisfied by
$x$ are the same as those encoded by $\mathcal{C}$, since admissibility
is a property of the physical system rather than of the training
distribution. Therefore $\pi^{-1}(\pi(x)) \in \Adm(X)$, and the
reconstruction is admissible. Distributional proximity is irrelevant;
admissibility membership is sufficient.
\end{proof}

The theorem makes precise what the \FSA\ zero-shot results demonstrate
empirically: it is the admissibility conditions, not the training distribution,
that determine the scope of generalization. A system that learns constraint
compatibility generalizes wherever those constraints apply. A system that
learns distributional statistics generalizes only where those statistics hold.
The practical consequence is that geometry-aware, constraint-preserving
architectures are qualitatively more robust than distribution-matching
architectures at any finite scale.

This reframes the scaling debate directly. The relevant question is not how
many parameters or how much data but whether the system has learned
the constraint structure of the domain. A small system that has learned
admissibility geometry may generalize more robustly than a large system
that has memorized distributional surface statistics. This is a strong
prediction that the current convergence of architectural pressure toward
field-theoretic design makes increasingly testable.

\subsection{Trajectory Coherence Precedes Semantic Interpretation}

A subtler but important observation emerges from both papers: coherent
dynamical organization appears to be detectable before and independently
of semantic interpretation. In the \FSA, the latent space organizes into
seasonally structured cyclic trajectories and historically directional drift
without any explicit semantic labelling. The geometric structure emerges
from the architecture's constraint commitments alone, prior to any
downstream task that would require semantic interpretation of the
compressed state.

In the moiré experiment, the coherent space-time propagation of the
collective modes is observed and characterized --- velocity, dispersion
relation, directionality, ballistic versus diffusive behaviour --- before the
exact nature of the IVC ground state is identified. The Goldstone mode is
recognizable as coherent propagation long before the question of whether
the underlying state is an incommensurate Kekul\'{e} spiral order or some
other IVC variant is resolved. Coherence geometry is a more primitive
experimental observable than phase classification.

This convergence suggests a broader epistemic principle: coherence
detection may be prior to semantic classification in the order of scientific
knowledge. Dynamical systems organize themselves coherently before human
observers interpret what the coherence means. The implication for machine
learning is that training coherence --- trajectory consistency across temporal
and spatial scales --- may be a more fundamental objective than semantic
correctness, and that systems evaluated primarily on semantic endpoint
accuracy may systematically miss the coherence structure that makes correct
endpoints possible at long horizons.

\subsection{Trajectory Coherence Under Compression}

The following theorem makes precise the sense in which admissibility-preserving
compression protects trajectory coherence, and characterises what happens
when it fails.

\begin{theorem}[Trajectory Coherence Preservation]
\label{thm:coherence}
Let $\pi: X \to M$ be a compression map that preserves admissibility and
cross-scale compatibility. If $\gamma: [t_1, t_2] \to X$ is a globally
admissible trajectory in $X$, then the projected trajectory
$\pi \circ \gamma$ preserves coherence under temporal composition:
\[
(\pi \circ \gamma)(t_2) \circ (\pi \circ \gamma)(t_1)
= \pi\bigl(\gamma(t_2) \circ \gamma(t_1)\bigr),
\]
up to bounded reconstruction error determined by the local entropy density
$S_t$. Conversely, if $\pi$ fails to preserve admissibility structure,
temporal composition becomes non-associative under projection, producing
cumulative trajectory drift and eventual \emph{representational
delamination}: the projected trajectory separates from the physical
admissibility manifold and cannot be returned to it by any sequence of
admissible local corrections.
\end{theorem}

\begin{proof}[Sketch]
Admissibility-preservation implies that $\pi$ is injective on $\Adm(X)$,
so the composition law in $X$ transfers to $M$: $\pi(\gamma(t_2) \circ
\gamma(t_1)) = (\pi \circ \gamma)(t_2) \circ (\pi \circ \gamma)(t_1)$
with error bounded by the reconstruction error of $\pi$, which is
controlled by $S_t$ (Proposition~\ref{prop:sparse}). If $\pi$ is not
admissibility-preserving, $\pi$ collapses some admissible transitions:
there exist $\omega_1 \neq \omega_2 \in \Adm(X)$ with
$\pi(\omega_1) = \pi(\omega_2)$. Temporal composition applied to such
collapsed transitions produces a sequence of projected states that
satisfies the projected composition law but not the original one; the
error accumulates over the trajectory and is not bounded by any local
quantity. Topological delamination follows once the accumulated error
exceeds the admissibility threshold $\kappa_*$.
\end{proof}

The phrase \textit{representational delamination} deserves emphasis. It
names the specific failure mode that distinguishes structurally
admissibility-violating systems from merely noisy ones. A noisy system
produces small random errors that may be correctable locally. A delaminated
system has left the physical admissibility manifold entirely; its trajectory
in representation space no longer corresponds to any trajectory in the
physical system's state space, and no amount of local correction can
restore the correspondence. Climate trajectory drift that exits the physical
climate manifold, latent space collapse that loses the seasonal trajectory
structure, semantic incoherence that accumulates over long text generation,
and hallucination that produces locally plausible but globally inadmissible
outputs are all instances of representational delamination. The \FSA\
architecture resists delamination by construction; the convolutional
baseline exhibits the precursor to it in its zero-shot resolution experiments.

\subsection{Constraint Curvature and Historical Drift}

The latent warming drift observed in the \FSA\ representations is not merely
a visualisation artifact. It suggests that the admissibility manifold
possesses an intrinsic curvature induced by historical constraint
accumulation. The climate trajectory does not translate uniformly through
latent space; it follows a curved path whose geometry is history-dependent
in a way that distinguishes accessible future states from counterfactually
accessible past states.

In the \RSVP\ framework this curvature arises from the progressive
contraction of the accessibility structure:
\[
\A_{t+1}^{\mathrm{consistent}} \subsetneq \A_t^{\mathrm{counterfactual}}.
\]
The resulting geometry is non-Euclidean in a precise sense. Parallel
transport of trajectories through the admissibility manifold produces
path-dependent displacement analogous to holonomy in curved connections.
Two trajectories that begin arbitrarily close may diverge permanently if
their histories induce different admissibility contractions --- the
topological manifestation of irreversibility discussed in \S\ref{sec:novel}.

Constraint curvature therefore provides a unified geometric interpretation
of four phenomena that otherwise appear as separate empirical observations:
the directional warming drift in the \FSA\ latent space (curvature in the
temporal direction of the accessibility manifold); the path-dependent
propagation of the spin--valley Goldstone modes (curvature in the spatial
direction of the constraint connection); the non-vanishing holonomy of the
sheaf formalism in \S\ref{sec:sheaf} (topological curvature of the
admissibility bundle); and the long-horizon coherence failures of systems
that ignore constraint geometry (geodesic deviation in a curved manifold
when a flat approximation is assumed). The latent drift is not a
visualisation curiosity. It is a geometric record of the curvature that
the physical system has accumulated.

\subsection{The Structural Pressure Hypothesis}

The formal correspondences established in the preceding sections are not,
by themselves, evidence that the \RSVP/\CLIO/\TARTAN\ framework is correct.
What makes them significant is their character and their independence. The
\FSA\ paper was designed to solve a specific engineering problem: scalable
climate emulation. The spin--valley moiré paper was designed to solve a
specific measurement problem: directly imaging hidden collective dynamics in
a correlated quantum material. The \RSVP/\CLIO/\TARTAN\ framework was
developed to address foundational questions in the philosophy of physics and
computation. All three arrived at overlapping structural commitments through
entirely different channels.

The \FSA\ demonstrates computationally that effective compression requires
geometry-aware, constraint-preserving, hierarchically explicit architecture.
The moiré experiment demonstrates physically that effective understanding of
collective dynamics requires sensitivity to hidden constraint manifolds that
are invisible to standard observable projections. Together they constitute
two independent instances of the same structural principle: the
dynamically significant organization of a physical system is carried by its
constraint geometry, and observable outputs are secondary projections of
that geometry rather than primary objects of study.

We call this the \textit{structural pressure hypothesis}: geometry-aware,
constraint-preserving, trajectory-sensitive representations emerge
independently wherever the physical demands of the target domain make
observable-surface approaches fail. The \FSA\ paper is one instance; the
moiré experiment is another; and similar pressures are visible in
developments in protein structure prediction~\cite{Lam2023,Bi2023}, neural
operator theory, and long-horizon planning in reinforcement learning. The
convergence is not coincidental. It reflects a structural feature of
physical reality: the dynamically significant organization of a system is
carried by its constraint geometry, not by its instantaneous observable
configuration.

\subsection{What the Convergence Does Not Imply}

It is equally important to be explicit about what the convergence does not
imply. It does not imply that the \FSA\ architecture is a complete
engineering implementation of the \RSVP\ programme, or that the spin--valley
experiment constitutes a laboratory demonstration of \RSVP\ field dynamics.
Both empirical papers operate within their own conceptual frameworks and
make no reference to the theoretical structures developed here. The
correspondences are structural, not definitional.

More importantly, the convergence does not imply that the engineering and
experimental approaches are approaching a theoretical limit that the formal
framework has already reached. Both the theoretical and empirical traditions
are works in progress. The theoretical framework identifies structural
features that the empirical approaches lack --- explicit admissibility
formalisms, provenance structures, cohomological constraint checking,
holonomy-aware path conditioning. The empirical approaches provide
constraints and exemplars that the theoretical framework has not yet been
tested against in those specific physical regimes. The relationship is one
of productive tension rather than one-sided validation.

\subsection{Long-Horizon Coherence as the Common Constraint}

The structural pressure hypothesis identifies a specific common constraint
that unifies the climate system, the spin--valley moiré experiment, and the
broader domains of cognition, reasoning, and semantic memory: \textit{local
correctness is insufficient for global admissibility}. This is the deep
unifying principle beneath the essay's entire argument, and it deserves
explicit statement as a subsection rather than a remark distributed across
the text.

In the climate system, a model that is locally accurate --- that predicts
the correct atmospheric state at each individual time step --- can
nonetheless produce a globally inadmissible trajectory, one that drifts to
climatologically impossible states over long integration horizons. The
local correctness of each step provides no guarantee that the global
trajectory remains on the physical manifold. Enforcing global admissibility
requires a representational structure that is sensitive to the full
constraint geometry of the system's possible trajectories, not merely to
the local statistics of adjacent states.

In the spin--valley system, a measurement strategy that is locally accurate
--- that correctly records the charge density at each spatial point --- can
nonetheless be entirely blind to the globally organized collective modes
that carry the system's physical order. The local measurement correctness
provides no access to the global constraint structure that the Goldstone
modes propagate. Accessing that structure requires a measurement approach
that is sensitive to the system's trajectory in hidden field space, not
merely to its projection onto the observable surface.

In cognition and language, a generative system that is locally accurate ---
that produces grammatically and semantically plausible continuations at each
step --- can nonetheless produce globally incoherent outputs: texts that
are locally reasonable but collectively contradictory, plans that are
individually achievable but jointly infeasible, arguments that are locally
valid but globally fallacious. Enforcing global coherence requires precisely
the trajectory-conditioning and provenance-tracking mechanisms that the
\RSVP/\CLIO\ framework formalizes.

The structural unity of these cases is not coincidental. In each, the system
has a long-horizon constraint structure that is not reducible to any finite
set of local constraints. Short-horizon optimization fails not because the
local objectives are wrong but because global admissibility is a
\textit{path-dependent} condition: whether a current state is admissible
depends on the entire history of transitions that produced it, not merely
on its current configuration. This is the \RSVP\ account of accessibility
shrinkage made precise as an engineering and experimental constraint.
Systems that do not maintain provenance-sensitive access to the full
history of admissible transitions cannot enforce global admissibility,
regardless of their local accuracy.

\subsection{The Remaining Conceptual Gap}

The conceptual gap that separates the current engineering approaches from
the full theoretical framework can be stated precisely. Existing
architectures, including the \FSA, achieve operational admissibility
preservation: their outputs tend to satisfy physical admissibility
conditions because the training data is organized around such conditions.
They do not achieve structural admissibility preservation: they lack
explicit formal mechanisms for enforcing admissibility in the event of
distribution shift, novel boundary conditions, or compositionally novel
inputs.

The distinction matters because it determines generalization. An
architecture that achieves only operational admissibility preservation will
generalize within the training distribution and fail at its boundary. An
architecture that achieves structural admissibility preservation will
generalize to any input that satisfies the same admissibility conditions as
the training data, regardless of distributional proximity. The \FSA's
zero-shot super-resolution capability is a partial demonstration of
structural generalization, which is why it is theoretically significant.
But it remains partial: the architecture has learned constraint
compatibility within a specific domain and a specific distribution, not
the general admissibility structure of the underlying physics.

The full theoretical framework provides a research direction for bridging
this gap. Architectures equipped with explicit admissibility formalisms,
sheaf-theoretic constraint checkers, and holonomy-aware trajectory
conditioning would, in principle, achieve structural rather than merely
operational admissibility preservation. Building and evaluating such
architectures is a concrete research programme that the \FSA\ paper helps
to motivate.

\subsection{Why Endpoint Metrics Systematically Mislead}

Most contemporary evaluation protocols assess endpoint states rather than
admissible trajectories. Reconstruction RMSE, next-token prediction loss,
classification accuracy, and benchmark pass rates all compare final outputs
against target outputs while remaining largely insensitive to the constraint
geometry that produced them. These metrics are degenerate observables in the
sense of \S\ref{sec:rsvp}: they collapse structurally distinct systems into
the same observable equivalence class.

The \FSA\ paper demonstrates this directly and precisely. The HEALPix
convolution baseline achieves competitive RMSE at low compression ratios
--- it is not a poor model by the standard metric. Yet its latent space is
disorganised, its zero-shot generalisation fails, and its parameter scaling
exhibits instability (validation loss diverging as parameters increase beyond
36M). The Field-Space architecture achieves lower RMSE at higher compression
while simultaneously producing organised latent trajectories and stable
parameter scaling. Both dimensions of improvement are missed by endpoint
RMSE alone.

The asymmetry is structural: endpoint metrics reward local correctness
without penalising global inadmissibility. A model that accurately
reconstructs each individual state while violating the dynamical geometry
of the trajectory between states scores well. A model that preserves
trajectory geometry while differing superficially from the training
distribution may score poorly. Benchmark culture therefore systematically
selects for representational systems that optimise projection-surface
fidelity at the expense of hidden constraint coherence.

The appropriate evaluation metric for an admissibility-preserving system is
not endpoint RMSE but the quality of the learned constraint structure,
measured by properties such as latent space organisation, zero-shot
generalisation to novel admissibility conditions, and long-horizon
trajectory coherence. These metrics are harder to compute but correspond
more directly to what the physical system actually requires.

\subsection{Broader Implications for Representation Theory}

The convergence also has implications beyond the specific domains of climate
modelling and field theory. If the structural pressure hypothesis is correct,
then the ongoing failure of large flat models to achieve long-horizon
coherence in language, planning, and reasoning is not primarily a scaling
problem but a structural one. The representational ontology of flat token
sequences, like the representational ontology of flat climate grids, is
incompatible with the geometry of the processes it is trying to model.
Increasing the scale of a fundamentally mismatched representation does not
resolve the mismatch; it may, in some cases, amplify it by increasing the
number of degrees of freedom available for overfitting to surface
regularities~\cite{Yan2024,Hunt2025}.

A direct consequence concerns benchmark culture. Most evaluations assess
endpoint accuracy --- reconstruction RMSE, classification accuracy, next-token
prediction loss --- without examining admissibility-preserving trajectory
organization. The \FSA\ paper contains an implicit critique of this practice:
the HEALPix convolution baseline achieves competitive RMSE at low
compression ratios, yet its latent space is disorganized and its zero-shot
generalization fails. A model can appear accurate while structurally
violating the geometry of the underlying process. Surface metrics reward
local correctness without penalizing global inadmissibility. The appropriate
evaluation metric for an admissibility-preserving system is not endpoint
RMSE but the quality of the learned constraint structure --- measured by
properties such as latent space organization, zero-shot generalization to
novel admissibility conditions, and long-horizon trajectory coherence.

The geometry-first, constraint-preserving, hierarchically explicit approach
suggests a different design and evaluation strategy: rather than increasing
the depth and width of flat representations, increase the explicitness and
precision of the constraint structure, and evaluate on the properties that
depend on that structure. This is not a proposal to abandon statistical
learning. It is a proposal to provide statistical learning with a
representational substrate that is geometrically compatible with the target
domain, so that what is learned reflects the admissibility structure of the
domain rather than the distributional structure of the training corpus.

\subsection{Representation Collapse as a General Failure Mode}

The recurring failure identified across climate modelling, quantum
measurement, and large-scale machine learning is not merely loss of accuracy
but \textit{collapse of representational fidelity to the underlying
constraint geometry}. A representation collapses when it preserves local
surface statistics while destroying the admissibility structure that
organises the physical process globally. The collapse is not always visible
in endpoint metrics because the metrics themselves are evaluated on the
projection surface rather than on the hidden trajectory geometry of the
process.

Representation collapse has a precise formal signature. A compressed
representation $M$ has undergone constraint collapse relative to the
original space $X$ if
\[
\dim\bigl(\A(M)\bigr) \ll \dim\bigl(\pi(\A(X))\bigr),
\]
where $\A(M)$ is the accessibility structure of the compressed
representation and $\pi(\A(X))$ is the image of the original accessibility
structure under compression. The compression has destroyed the accessible
future structure of the physical system: states that should be reachable
in the compressed representation are not, and states that should not be
reachable are.

This failure mode explains why systems can appear operationally successful
under short-horizon evaluation while remaining structurally incapable of
long-horizon coherence. Over short horizons, the collapsed accessibility
structure is approximately consistent with the original; the errors are small
enough to fall within acceptable metric tolerances. Over long horizons,
the divergence between $\A(M)$ and $\pi(\A(X))$ accumulates, and the
system eventually generates trajectories that are inadmissible in the
original space. This is representational delamination in slow motion.

The \FSA\ results provide a controlled empirical demonstration. The
convolutional baseline at high compression ratios undergoes partial
constraint collapse: it loses the organised accessibility structure visible
in the Field-Space latent trajectories. The short-horizon RMSE metric does
not detect this collapse until the extreme compression regime (1024$\times$)
where accuracy degrades sharply. The latent space organisation metric
detects it at every compression ratio.

% =============================================================
\section{Physics Education Research as a Third Instantiation}
\label{sec:edu}

A third domain provides unexpected confirmation of the structural pressure
hypothesis from a direction entirely orthogonal to both climate emulation
and condensed matter physics. A recent comparative study of active learning
methods in introductory physics and astronomy courses~\cite{Sundstrom2025}
examined four established instructional approaches across 31 implementations
at 28 institutions: Peer Instruction~\cite{Mazur1997}, ISLE~\cite{Etkina2007},
Tutorials~\cite{McDermott2002}, and SCALE-UP~\cite{Beichner2007}. The study
reported that all four methods improved short-term conceptual understanding
as measured by concept inventories, with SCALE-UP exhibiting significantly
larger gains. Its observational analysis suggested that time spent in genuine
collaborative work, rather than in lecture-plus-clicker cycles, predicted
larger conceptual gains.

This is a well-executed study within its design assumptions. It is precisely
those assumptions that make the paper theoretically interesting in the present
context. The study implicitly treats physics learning effectiveness as a
scalar quantity measurable through short-term standardized instruments. The
argument of this essay predicts, on structural grounds, that this compression
should produce exactly the failure mode described in preceding sections:
representational collapse of the underlying admissibility structure, invisibility
of the collapse to the metric used, and systematic misleading of long-horizon
inference from short-horizon evaluation.

\subsection{Concept Inventories as Degenerate Observables}

A concept inventory is a multiple-choice instrument designed to assess
conceptual understanding of a particular domain, typically administered
before and after an instructional intervention. Force Concept Inventories,
conceptual surveys of electricity and magnetism, and similar instruments have
been widely validated for reliability and discriminant validity within their
intended short-term measurement context. The question at issue is not their
reliability as instruments but their adequacy as observables of the underlying
process they are meant to track.

In the formal vocabulary developed in \S\ref{sec:rsvp}, a concept inventory
score is a degenerate observable of intellectual development. The projection
\[
\mathcal{O}_{\mathrm{inv}}: X_{\mathrm{dev}} \to \mathbb{R}
\]
maps the full developmental state $X_{\mathrm{dev}}$ of a learner --- which
includes conceptual structure, tacit knowledge~\cite{Polanyi1966},
problem-solving heuristics, analogical repertoire, motivational trajectory,
and the accumulated pattern of domain-crossing experience --- onto a single
scalar score. The map is degenerate in a structurally precise sense: it
collapses dynamically distinct developmental configurations into the same
observable value.

Two students who achieve identical pre-to-post gains on a force concept
inventory may have undergone radically different developmental trajectories.
One may have consolidated a robust, transferable conceptual schema with
rich connections to adjacent domains. Another may have memorised the
vocabulary sufficient to pass the instrument without acquiring the underlying
structural understanding. Both map to the same observable, but their
accessible futures --- the set of admissible developmental trajectories
available to them at the next stage --- are entirely different. The degenerate
observable $\mathcal{O}_{\mathrm{inv}}$ cannot distinguish between them.

This degeneracy is not a correctable measurement error that better instrument
design could eliminate. It is structural: the dimensionality of
$X_{\mathrm{dev}}$ far exceeds the dimensionality of $\mathbb{R}$, and no
scalar projection can avoid the collapse. The study by Sundstrom et al.\ is
careful and rigorous within its design, but it is asking about the image of
$\mathcal{O}_{\mathrm{inv}}$ rather than about $X_{\mathrm{dev}}$ itself.
The structural pressure hypothesis predicts that this will produce
operationally useful but long-horizon-misleading conclusions: the gains are
real, the measurement is coherent, and the inference to long-term scientific
capability is structurally unwarranted.

\subsection{The Reflexivity Complication}

The educational domain compounds the degeneracy problem with a second
structural feature absent from climate and condensed matter physics: reflexivity.
Physical fields do not modify their behaviour in response to being measured.
Climate systems do not read papers about climate modelling and adjust their
dynamics accordingly. Spin--valley modes do not reorganise themselves in
response to observational protocols.

Human learners do all of these things. Students respond to grading structures,
institutional incentives, perceived performance standards, and the cultural
meanings attached to different disciplines. A concept inventory is not merely
an instrument applied to a stable system; it is an intervention that modifies
the system it measures. Students who know they are being evaluated on
conceptual understanding may allocate effort differently than those who are
not. Instructors who know their effectiveness is being measured by pre-to-post
gains on a specific instrument may align their teaching toward that instrument.
The act of measurement modifies the accessibility structure of the measured
system in ways that have no analogue in physical measurement.

This reflexivity problem is formally analogous to the difference between
a constraint that is merely statistical and one that is structurally enforced.
A statistical constraint can be violated by a sufficiently motivated agent;
a structural constraint cannot be violated by any continuous adjustment. Human
learning systems operate under statistical constraints that bend under
institutional pressure. The result is Goodhart's Law~\cite{Campbell1979}:
when a measure becomes a target, it ceases to be a good measure. The
concept inventory score, once used to evaluate instructional effectiveness,
partially decouples from the construct it was designed to track, because the
measurement itself becomes part of the constraint structure that shapes
instructional decisions.

The RSVP framework handles reflexive systems by tracking the provenance
of the accessibility structure: $\Pp_t = \Pp(\omega_t)$ records not merely
the current state but the history of how the admissibility conditions
themselves have evolved. A non-reflexive physical system has a stable
admissibility structure that the measurement does not perturb. A reflexive
social system has an admissibility structure that is partially constituted
by the measurement protocol. Without explicit provenance tracking, an
analysis of a reflexive system cannot distinguish between genuine improvement
in the underlying process and improvement in the system's capacity to optimise
for the measurement.

\subsection{Intellectual Development as Admissibility Geometry}

The most theoretically productive framing of the education paper's results
is as evidence about accessible futures rather than about scalar quantities.
What SCALE-UP produces, if the observational analysis is correct, is an
instructional environment that expands the accessible developmental trajectory
space more effectively than the alternatives. Students who spend more time
in genuine collaborative work on complex problems develop a richer set of
admissible subsequent trajectories: they have more ways to proceed, more
analogical connections available, more problem-solving strategies in their
repertoire. The concept inventory gain is a degenerate projection of this
expanded accessibility; it is real evidence of the underlying change but not
a measure of it.

This reframing dissolves an apparent tension in the literature. Studies
occasionally find that active learning methods improve standardized test
scores without producing equivalent improvements in other measures of
scientific capability~\cite{Andrews2011}. The RSVP interpretation explains
this: the accessibility structure expanded by good instruction is
multidimensional, and a scalar projection will detect expansion in some
dimensions while being blind to others. Different instruments project
different dimensions; the apparent inconsistency across instruments reflects
genuine multidimensionality in the underlying developmental space rather than
measurement error.

The implications for educational research design are directly analogous to
the implications for computational and experimental design drawn from the
FSA and moiré papers. Just as effective climate compression requires
operators sensitive to the admissibility geometry of the atmospheric system
rather than to pixel-level statistics, and just as effective measurement
of spin--valley collective modes requires space-time imaging of hidden
trajectories rather than static charge response, effective assessment of
intellectual development requires instruments sensitive to the accessibility
structure of the learner's developmental space rather than to the endpoint
coordinates of a standardized projection.

Concretely: rather than measuring pre-to-post gains on a fixed instrument,
assess the structure of the learner's analogical transfer capacity across
novel domains, the coherence of their problem decomposition strategies over
long time horizons, and the stability of their conceptual representations
under perturbation by unusual examples. These are not merely harder
measurements; they are measurements of a different object --- the constraint
geometry of the developmental space rather than a degenerate scalar
projection of it. The analogy with the FSA's latent space organisation
metric versus its RMSE metric is direct: the latter is easier to compute
and partially informative; the former is more informative about the
structural property that determines long-horizon performance.

% =============================================================
\section{Constraint Geometry and the Failure of Flat Epistemology}
\label{sec:flat}

\subsection{Why Surface Statistics Become Misleading}

The recurring pattern across every empirical domain examined in this essay
is that systems initially appear tractable through surface statistics alone,
and then progressively fail as the temporal horizon, geometric complexity,
or coupling structure of the domain increases. Flat representations work
well in the weak-coupling regime because local statistics partially shadow
the deeper admissibility geometry. They fail in the strong-coupling regime
because the local statistics become progressively decoupled from the
constraint structure that actually governs long-horizon evolution.

This transition can be formalised geometrically. Let $X$ denote the full
constraint space of the physical system and let $\pi: X \to M$ be a
projection onto an observable manifold $M$. In weakly constrained systems,
the fibers $\pi^{-1}(m)$ for $m \in M$ remain small relative to the scale
of the dynamics. Observable states therefore approximately determine the
underlying physical state, and surface statistics are informative. In
strongly constrained systems, the fibers become large and structurally
heterogeneous: many dynamically distinct trajectories project to the same
observable state. The observable surface ceases to determine the physical
trajectory.

The result is a characteristic epistemic illusion. The system appears
predictable over short horizons because the observable projection still
contains enough information to interpolate locally. Over longer horizons,
prediction degrades because the hidden admissibility geometry dominates
the dynamics while remaining invisible to the representation. This is the
failure mode visible in long-context language generation, in climate
trajectory drift, in hidden quantum collective modes, and in educational
evaluation. Local correctness survives while global coherence collapses.

\subsection{Constraint Geometry as the Carrier of Persistence}

The deeper implication is that persistence itself is a geometric property.
A system persists through time not because its local state is stable but
because its admissibility structure constrains the space of possible
perturbations tightly enough that coherent trajectories remain available.

This suggests a reinterpretation of stability theory. Traditional dynamical
systems theory treats stability in terms of local perturbative return: a
trajectory is stable if nearby trajectories converge back toward it under
evolution. The present framework suggests a stronger criterion. A trajectory
is structurally stable if perturbations preserve admissibility membership:
$x + \delta x \in \Adm(X)$ for all perturbations $\delta x$ within some
admissible neighbourhood.

The distinction matters because a trajectory may be locally stable while
globally inadmissible. Such trajectories are operationally persistent over
short horizons but eventually undergo representational delamination from
the physical manifold. This is exactly what occurs in systems optimised
exclusively for local loss functions without explicit admissibility
preservation. The hidden field is therefore not merely explanatory overhead.
It is the carrier of persistence itself.

\subsection{Constraint Transport and the Nature of Generalisation}

The standard interpretation of generalisation in machine learning is that
models interpolate statistical regularities beyond the training set. The
evidence examined in this essay suggests a more precise interpretation:
systems generalise robustly when they learn transport laws for admissibility
constraints rather than surface correlations.

This can be stated formally. Let $\mathcal{C}(x)$ denote the constraint
structure associated with a state $x$. A system generalises structurally if
$\mathcal{C}(x_{\mathrm{train}}) = \mathcal{C}(x_{\mathrm{test}})$ implies
successful reconstruction or prediction even when $x_{\mathrm{test}} \notin
\mathcal{D}_{\mathrm{train}}$. This is exactly what the \FSA\ transformer
demonstrates in its transfer from ERA5 fields to MPI-ESM outputs: the
surface statistics differ, the constraint structure does not, and the
architecture transfers.

The implication reframes the scaling problem. Increasing parameter count
without improving constraint transport primarily increases interpolation
capacity inside the observable manifold. Improving admissibility preservation
increases the volume of structurally reachable trajectories outside the
training distribution. The former scales memory. The latter scales
understanding.

\subsection{Scientific Revolutions as Representational Topology Change}

The argument developed throughout this essay provides a structural analogy
--- offered not as a causal explanation but as a formal parallel --- for
why scientific paradigms shift. A representational framework remains
productive as long as its observable manifold remains approximately compatible
with the admissibility structure of the domain. As empirical precision
increases, the mismatch accumulates. Anomalies appear because trajectories
that are admissible in the representation are inadmissible in the physical
system. Eventually the representation becomes topologically unstable relative
to the domain, and adding parameters, corrections, or local refinements no
longer restores coherence. The representational substrate itself must change.

This structural description applies, with appropriate domain-specific
differences, to the shift from Euclidean to relativistic spacetime, from
local hidden-variable pictures to Hilbert-space quantum mechanics, from
symbolic AI to geometry-aware latent dynamics, and from flat climate grids
to field-space representations. In each case the observable surface
eventually ceases to support the admissible trajectories required by the
physics.

The structural parallel is this: scientific revolutions share with these
engineering transitions the feature that they cannot be resolved by
incremental local improvement. They require a change of representational
topology --- a change in the class of objects the framework treats as
primary, not merely a change in the parameters of a fixed framework.

\subsection{Toward a Constraint-First Epistemology}

Traditional empiricism treats observables as primary and hidden structure
as derivative inference. The repeated convergence examined in this essay
suggests the reverse ordering may be more fundamental: constraint geometry
is primary, while observables are projections that preserve only part of
the underlying organizational structure.

Knowledge, under this view, is not fundamentally the accumulation of surface
measurements. It is the progressive reconstruction of the hidden admissibility
manifold from the traces left by observable projections. The aim of science
changes subtly but decisively: the goal is not merely to predict observations
but to recover the constraint geometry that makes those observations cohere
over long horizons.

This explains why long-horizon coherence repeatedly emerges as the deepest
criterion of successful representation. A representation that captures the
constraint geometry continues to produce admissible trajectories under
perturbation, extrapolation, and scale transfer. A representation that
captures only surface regularities fails precisely when the system leaves
the narrow manifold occupied by the training distribution. The difference
between the two is the difference between prediction and understanding.

% =============================================================
\section{Future Directions}
\label{sec:future}

The convergence documented throughout this essay suggests several concrete
research directions whose importance extends beyond the specific domains
examined here.

\textit{Explicit admissibility-aware architectures.} Existing
geometry-aware systems such as the \FSA\ preserve admissibility
operationally but not formally. Architectures equipped with explicit
constraint-checking operators, sheaf-theoretic compatibility verification,
and trajectory-level coherence regularization would allow admissibility to
be enforced structurally rather than statistically. Such systems would
constitute a qualitatively different category of representation learner:
not distribution-matching systems but admissibility-preserving systems,
whose generalization failures correspond to genuine constraint violations
rather than distributional novelty.

\textit{Observable degeneracy detection.} Many scientific domains likely
contain hidden dynamical manifolds analogous to the spin--valley modes in
twisted WSe$_2$: physically organizing structures that standard observables
systematically project away. A general theory of projection degeneracy would
provide criteria for identifying when an observable basis is insufficient
to recover the underlying constraint geometry. The formal tools are already
present in \S\ref{sec:rsvp} and the sheaf formalism of \S\ref{sec:sheaf};
what remains is their systematic application to domain selection.

\textit{Developmental and cognitive trajectory analysis.} If educational and
cognitive systems possess rich accessibility structures, then scalar endpoint
evaluations systematically underestimate the geometry of learning.
Longitudinal trajectory-sensitive metrics, analogical transfer topology, and
provenance-aware developmental analysis may become necessary for understanding
cognition in the same way that space-time imaging became necessary for
understanding the spin--valley superfluid. The Sundstrom et al.\ study
measures concept inventory gains; the analogous advancement would measure
the structure of the learner's accessible trajectory space directly.

\textit{Latent geometry as an attractor under optimization.} The spontaneous
emergence of organised latent manifolds in \FSA\ systems trained only for
reconstruction suggests that certain geometric structures are attractors
under optimization pressure whenever the underlying domain contains persistent
admissibility constraints. Understanding this relationship would help explain
why increasingly many high-performing systems converge toward field-like
internal organizations even when not explicitly designed to do so, and would
provide principled guidance for when such convergence should be expected to
occur.

\textit{Unification across projection types.} The present essay has treated
physical compression, quantum measurement, and educational assessment as
three instances of the same structural problem. If constraint geometry is
the primary organizational substrate across these domains, then field
dynamics, inference, memory, and measurement may be different projections
of a common underlying structure. That possibility remains speculative. The
structural convergence motivating it does not.

% =============================================================
\section{Conclusion: The Hidden Field is the Physical Field}
\label{sec:conclusion}

Three empirical papers, from three entirely different domains, each
independently discover that the dynamically significant organization of their
respective systems is not carried by the obvious observable surface. The
Field-Space Autoencoder~\cite{Meuer2026} finds that climate compression
fails when it optimises for pixel-level reconstruction fidelity, and succeeds
when it preserves the admissibility geometry of the atmospheric field. The
spin--valley moiré experiment~\cite{Xiong2026} finds that the physically
meaningful collective modes of the quantum material are charge-decoupled and
invisible to standard observational bases, becoming accessible only through
direct space-time imaging of hidden field trajectories. The physics education
study~\cite{Sundstrom2025} measures short-horizon gains on scalar concept
inventories, instruments that are degenerate observables of the multidimensional
accessibility structure of intellectual development, producing results that are
real but structurally incomplete as accounts of long-horizon scientific
capability.

The three cases form a progression from the most physical to the most reflexive,
and the structural argument scales across all three. In the climate system,
the hidden constraint geometry is the admissibility structure of the
atmospheric field, which the FSA learns to compress without destroying. In
the quantum material, the hidden constraint geometry is the spin--valley order
parameter and its collective modes, which the pump--probe imaging makes
directly accessible. In the educational system, the hidden constraint geometry
is the multidimensional developmental accessibility structure of the learner,
which no scalar instrument can capture and which reflexivity partially
constitutes through the measurement process itself.

This progression reveals a gradient in the severity of the representational
problem. Physical systems have stable admissibility structures that are in
principle fully accessible to appropriate measurement strategies. The problem
is choosing the right measurement basis. Human developmental systems have
admissibility structures that are historically contingent, reflexively
modified by measurement, and entangled with the entire biographical and
institutional context of the learner's development. The problem is not merely
choosing the right measurement basis but acknowledging that no finite set of
standardized measurements can capture the full constraint geometry of a
historically embedded developmental trajectory.

The convergence has three implications, at increasing levels of generality.

The first is methodological. Standard observational strategies across all
three domains are systematically biased toward the observable projection
surface. Climate models that optimize for pixel-accurate reconstruction,
condensed matter measurements that record only charge response, and
educational studies that assess short-term standardized gains are all
measuring their respective projections rather than the underlying processes.
The structural pressure hypothesis predicts that this bias will produce
systematic failures at precisely the scales and regimes where the constraint
geometry diverges most strongly from the observable projection. These failures
are not evidence of insufficient data or insufficient parameters; they are
evidence of a representational mismatch.

The second is theoretical. The \RSVP\ framework provides a vocabulary in
which all three findings can be stated with uniform precision. Constraint-preserving
projection, admissibility-preserving compression, degenerate observables,
accessible future structure, representational delamination, and topological
obstruction apply equally to atmospheric dynamics, collective quantum modes,
and intellectual development. The formal vocabulary transfers because the
structural problem is the same: how to represent a system with rich
constraint geometry using limited observational resources without destroying
the geometric structure that makes long-horizon behaviour predictable.

The third is philosophical. The repeated independent discovery of constraint
geometry as the carrier of physical and developmental organization --- across
physical, quantum, and human domains --- suggests that this is not a feature
of particular theoretical frameworks but a structural feature of reality as
systems that persist and organize themselves. Physical systems organize
along admissible trajectories through constraint manifolds. The macroscopic
collective behaviour we observe is the projection of that hidden organization
onto the observational surfaces we have developed to detect it. The surfaces
are real, but they are not the primary objects. The primary objects are
the constraint manifolds and the trajectories through them.

The geometry is not optional. It is what the physics is made of.

% =============================================================

\vspace{2em}
\noindent\rule{\linewidth}{0.4pt}
\vspace{0.5em}

{\small\noindent\textbf{Acknowledgement.} This work was developed as an
independent research contribution without institutional affiliation.
All formal frameworks --- \RSVP, \CLIO, \TARTAN\ --- have been developed
under the name Flyxion. The present essay situates them in relation to
concurrent developments in scientific machine learning, condensed matter
physics, and physics education research, none of which was consulted in
the development of the frameworks.}

\vspace{2em}

% =============================================================
\appendix

\section{Category-Theoretic Formalization}
\label{app:category}

The formal structures deployed throughout this essay can be given a unified
category-theoretic interpretation. This appendix states the principal
objects and morphisms explicitly, for readers who prefer the categorical
perspective as a way of tracking the logical dependencies between the
framework's components.

\subsection*{The Trajectory Category}

Let $\mathbf{Traj}$ denote the category whose objects are admissible state
spaces $X$ equipped with an accessibility structure $\A_t$ and an
admissibility function $\kappa: \Omg \times H \to \mathbb{R}$, and whose
morphisms are admissibility-preserving maps:
\[
f \in \mathrm{Hom}_{\mathbf{Traj}}(X, Y)
\;\Longleftrightarrow\;
\omega \in \Adm(X) \Rightarrow f(\omega) \in \Adm(Y).
\]
Composition is the usual composition of maps, which preserves
admissibility by transitivity. The identity morphism on each object is the
identity map, which trivially preserves admissibility.

\subsection*{The Representation Functor}

Let $\mathbf{Rep}$ denote the category of representational manifolds:
compact smooth manifolds $M$ equipped with a metric and a latent
trajectory structure. The projection functor
\[
\Pi: \mathbf{Traj} \to \mathbf{Rep}
\]
maps each admissible state space $X$ to a compressed representational
manifold $M = \Pi(X)$ and each admissibility-preserving morphism to a
structure-preserving map between manifolds. The \FSA\ compression
operators $\mathcal{F}_{k \to j}$ are the component maps of $\Pi$ at
each scale stratum; the conditions they must satisfy to be
admissibility-preserving are precisely the cross-scale compatibility
conditions enforced by the field-space architecture.

A compression $\pi \in \mathrm{Hom}_{\mathbf{Rep}}(M, M')$ is
\textit{trajectory-preserving} if the following diagram commutes:
\[
\begin{tikzcd}[column sep=large, row sep=large]
X \arrow[r, "\Pi"] \arrow[d, "f"'] & M \arrow[d, "\pi"] \\
Y \arrow[r, "\Pi"'] & M'
\end{tikzcd}
\]
for every admissibility-preserving morphism $f: X \to Y$ in $\mathbf{Traj}$.

\subsection*{The Verification Natural Transformation}

The \CLIO\ verification map $\mathcal{V}: \widehat{\mathcal{L}} \to
\Svalid \cup \{\bot\}$ is a natural transformation between functors:
\[
\mathcal{V}: \mathcal{R} \circ \Pi \Rightarrow \mathrm{Id}_{\Svalid}
\cup \{\bot\},
\]
where $\mathcal{R}: \mathbf{Rep} \to \widehat{\mathbf{Rep}}$ is the
stochastic rendering functor. Naturality requires that verification
commutes with admissibility-preserving morphisms: if $f: X \to Y$ in
$\mathbf{Traj}$, then verifying the rendering of $\Pi(X)$ is consistent
with verifying the rendering of $\Pi(Y)$ after applying $f$.

\subsection*{The Accessibility Sheaf}

The accessibility sheaf $\scrA$ is a sheaf of sets over the model space
$\fM$, assigning to each open $U \subseteq \fM$ the set of admissible
transitions within $U$. Restriction maps are the canonical inclusions.
The gluing axiom requires that locally admissible transitions be globally
consistent: for any compatible family of admissible transitions on an open
cover $\{U_i\}$, there exists a unique globally admissible transition
restricting to each local one.

The \v{C}ech cohomology $\check{H}^\bullet(\mathfrak{U}, \scrA)$ measures
the failure of global sections to exist. In degree 1, a non-trivial class
$[\delta\sigma] \in \check{H}^1(\mathfrak{U}, \scrA)$ corresponds to a
globally incoherent transition system: the local admissibility conditions
are mutually incompatible in a topologically irresolvable way. The
full pipeline of \S\ref{sec:sheaf} fits into the long exact sequence of
cohomology associated to a short exact sequence of sheaves encoding the
inclusion of admissible transitions into all transitions.

\subsection*{Power and Holonomy}

The holonomy of the constraint connection $\Psi_t$ around a closed loop
$\gamma$ in $\fM$ defines an element of $\mathrm{Aut}(\Omg_t)$:
\[
\mathrm{Hol}_\gamma(\Psi) = \mathcal{P}\exp\!\Bigl(
\oint_\gamma A_\Psi\,d\ell\Bigr) \in \mathrm{Aut}(\Omg_t).
\]
Non-trivial holonomy implies that the admissibility structure is
path-dependent: the set of admissible transitions at a point in $\fM$
depends on the history of the path used to reach that point. This is
the categorical expression of the \RSVP\ irreversibility thesis:
provenance matters, and it cannot be recovered from instantaneous state.

The group of holonomies $\{\mathrm{Hol}_\gamma(\Psi) \mid \gamma
\text{ closed loop in } \fM\}$ is the holonomy group of the constraint
connection. Flat constraint regions (where $\mathfrak{R}_{ij} = 0$) have
trivial holonomy group; the historical ordering of transitions is
irrelevant. Non-flat regions have non-trivial holonomy; the historical
ordering is dynamically significant. The \FSA\ latent warming drift is
evidence of non-trivial holonomy in the climate system's constraint
connection over the accessible state space.

\begin{thebibliography}{99}

\bibitem{Meuer2026}
Meuer, J., Witte, M., Pl\'{e}siat, \'{E}., Ludwig, T.\ \& Kadow, C.
Field-space autoencoder for scalable climate emulators.
\textit{npj Artificial Intelligence} \textbf{2}, 50 (2026).

\bibitem{Xiong2026}
Xiong, R.\ \textit{et al.}
Observation of propagating collective spin--valley modes in twisted WSe$_2$.
\textit{Nature Physics} (2026).
\texttt{https://doi.org/10.1038/s41567-026-03280-w}

\bibitem{NatPhys2026SpinValley}
Taking snapshots of spin--valley modes in a moir\'{e} superlattice.
\textit{Nature Physics} (2026).
\texttt{https://doi.org/10.1038/s41567-026-03295-3}

\bibitem{BiFu2021}
Bi, Z.\ \& Fu, L.
Excitonic density wave and spin--valley superfluid in bilayer transition
metal dichalcogenide.
\textit{Nature Communications} \textbf{12}, 642 (2021).

\bibitem{MakShan2022}
Mak, K.~F.\ \& Shan, J.
Semiconductor moir\'{e} materials.
\textit{Nature Nanotechnology} \textbf{17}, 686--695 (2022).

\bibitem{Xia2025}
Xia, Y.\ \textit{et al.}
Superconductivity in twisted bilayer WSe$_2$.
\textit{Nature} \textbf{637}, 833--838 (2025).

\bibitem{Guo2025}
Guo, Y.\ \textit{et al.}
Superconductivity in $5.0^\circ$ twisted bilayer WSe$_2$.
\textit{Nature} \textbf{637}, 839--845 (2025).

\bibitem{Bernevig2025}
Bernevig, B.~A.\ \textit{et al.}
Fractional quantization in insulators from Hall to Chern.
\textit{Nature Physics} \textbf{21}, 1702--1713 (2025).

\bibitem{Witte2025}
Witte, M., Meuer, J., Pl\'{e}siat, \'{E}.\ \& Kadow, C.
Field-space attention for structure-preserving earth system transformers.
arXiv:2501.xxxxx (2025).

\bibitem{Lam2023}
Lam, R.\ \textit{et al.}
GraphCast: learning skillful medium-range global weather forecasting.
\textit{Science} \textbf{382}, 1416--1421 (2023).

\bibitem{Bi2023}
Bi, K.\ \textit{et al.}
Accurate medium-range global weather forecasting with 3D neural networks.
\textit{Nature} \textbf{619}, 533--538 (2023).

\bibitem{Kurth2023}
Kurth, T.\ \textit{et al.}
FourCastNet: accelerating global high-resolution weather forecasting using
adaptive Fourier neural operators.
\textit{Proceedings of the International Conference for High Performance
Computing, Networking, Storage and Analysis}, 1--14 (2023).

\bibitem{Rombach2022}
Rombach, R., Blattmann, A., Lorenz, D., Esser, P.\ \& Ommer, B.
High-resolution image synthesis with latent diffusion models.
\textit{Proceedings of CVPR}, 10684--10695 (2022).

\bibitem{Connor2021}
Connor, M., Canal, G.\ \& Rozell, C.
Variational autoencoder with learned latent structure.
\textit{Proceedings of AISTATS}, 2359--2367 (2021).

\bibitem{Gorski2005}
G\'{o}rski, K.~M.\ \textit{et al.}
HEALPix: a framework for high-resolution discretization and fast analysis of
data distributed on the sphere.
\textit{Astrophysical Journal} \textbf{622}, 759--771 (2005).

\bibitem{Hersbach2020}
Hersbach, H.\ \textit{et al.}
The ERA5 global reanalysis.
\textit{Quarterly Journal of the Royal Meteorological Society}
\textbf{146}, 1999--2049 (2020).

\bibitem{Song2021}
Song, J., Meng, C.\ \& Ermon, S.
Denoising diffusion implicit models.
\textit{International Conference on Learning Representations} (2021).

\bibitem{Lessig2023}
Lessig, C.\ \textit{et al.}
AtmoRep: a stochastic model of atmosphere dynamics using large scale
representation learning.
arXiv:2308.13280 (2023).

\bibitem{Friston2010}
Friston, K.
The free-energy principle: a unified brain theory?
\textit{Nature Reviews Neuroscience} \textbf{11}, 127--138 (2010).

\bibitem{Tononi2016}
Tononi, G., Boly, M., Massimini, M.\ \& Koch, C.
Integrated information theory: from consciousness to its physical substrate.
\textit{Nature Reviews Neuroscience} \textbf{17}, 450--461 (2016).

\bibitem{ShawStevens2025}
Shaw, T.~A.\ \& Stevens, B.
The other climate crisis.
\textit{Nature} \textbf{639}, 877--887 (2025).

\bibitem{Arora2018}
Arora, S., Hu, W.\ \& Kothari, P.~K.
An analysis of the t-SNE algorithm for data visualization.
\textit{Conference on Learning Theory}, 1455--1462 (2018).

\bibitem{Angiulli2022}
Angiulli, F., Fassetti, F.\ \& Ferragina, L.
LatentOut: an unsupervised deep anomaly detection approach exploiting latent
space distribution.
\textit{Machine Learning} \textbf{112}, 4323--4349 (2022).

\bibitem{Yan2024}
Yan, C.-W.\ \textit{et al.}
Fourier amplitude and correlation loss: beyond using L2 loss for skillful
precipitation nowcasting.
\textit{Advances in Neural Information Processing Systems} \textbf{37},
100007--100041 (2024).

\bibitem{Hunt2025}
Hunt, K.~M.~R.
Stop using root-mean-square error as a precipitation target!
arXiv:2509.08369 (2025).

\bibitem{Sundstrom2025}
Sundstrom, M., Gambrell, J., Green, C., Traxler, A.~L.\ \& Brewe, E.
Relative benefits of different active learning methods to conceptual physics
learning.
arXiv:2505.04577 [physics.ed-ph] (2025).

\bibitem{Mazur1997}
Mazur, E.
\textit{Peer Instruction: A User's Manual}.
Prentice Hall, Upper Saddle River, NJ (1997).

\bibitem{Etkina2007}
Etkina, E.\ \& Van Heuvelen, A.
Investigative science learning environment --- a science process approach to
learning physics.
In \textit{Research-Based Reform of University Physics}, vol.~1, pp.~1--48
(2007).

\bibitem{McDermott2002}
McDermott, L.~C.\ \& Shaffer, P.~S.
\textit{Tutorials in Introductory Physics}.
Prentice Hall, Upper Saddle River, NJ (2002).

\bibitem{Beichner2007}
Beichner, R.~J.\ \textit{et al.}
The student-centered activities for large enrollment undergraduate programs
(SCALE-UP) project.
In \textit{Research-Based Reform of University Physics}, vol.~1, pp.~2--39
(2007).

\bibitem{Andrews2011}
Andrews, T.~M., Leonard, M.~J., Colgrove, C.~A.\ \& Kalinowski, S.~T.
Active learning not associated with student learning in a random sample of
college biology courses.
\textit{CBE---Life Sciences Education} \textbf{10}, 394--405 (2011).

\bibitem{Campbell1979}
Campbell, D.~T.
Assessing the impact of planned social change.
\textit{Evaluation and Program Planning} \textbf{2}, 67--90 (1979).

\bibitem{Polanyi1966}
Polanyi, M.
\textit{The Tacit Dimension}.
Doubleday, Garden City, NY (1966).

\end{thebibliography}

\end{document}